% Standalone uncompiled TeX source. No PDF or layout verification is claimed.
% Companion verification files cover only their explicitly stated fragments.
\documentclass[11pt]{article}
\usepackage[T1]{fontenc}
\usepackage[utf8]{inputenc}
\usepackage[margin=1in]{geometry}
\usepackage{amsmath,amssymb,amsthm,mathtools,enumitem}
\usepackage[hidelinks]{hyperref}
\newtheorem{theorem}{Theorem}[section]
\newtheorem{proposition}[theorem]{Proposition}
\newtheorem{lemma}[theorem]{Lemma}
\newtheorem{corollary}[theorem]{Corollary}
\theoremstyle{definition}
\newtheorem{definition}[theorem]{Definition}
\newtheorem{example}[theorem]{Example}
\theoremstyle{remark}
\newtheorem{remark}[theorem]{Remark}
\newcommand{\R}{\mathbb R}
\newcommand{\N}{\mathbb N}
\newcommand{\Z}{\mathbb Z}
\newcommand{\id}{\operatorname{id}}
\newcommand{\im}{\operatorname{im}}
\newcommand{\supp}{\operatorname{supp}}
\newcommand{\Hom}{\operatorname{Hom}}
\newcommand{\Cap}{\operatorname{Cap}}
\newcommand{\sat}{\operatorname{sat}}
\newcommand{\cl}{\operatorname{cl}}
\newcommand{\Obj}{\operatorname{Ob}}
\newcommand{\Reach}{\operatorname{Reach}}
\newcommand{\dist}{\operatorname{dist}}
\title{Endogenous Capability Growth through Native Interaction\\
\large Manufacture, reusable interfaces, and subsequent experiments}
\author{}
\date{5 November 2026}
\begin{document}
\maketitle
\begin{abstract}
Interaction can manufacture an artifact whose warranted service was not
available in the previous installed configuration. This note develops
that claim with explicit native constructions, matched before-and-after
service comparisons, and limits on empirical warrant. The complete
manufacture-and-use protocol remains inside the supplied runtime; no
transcendence of its extensional closure is claimed.

Finite mass-conserving reaction systems admit exact manufacture and
service synthesis, including a bounded cost-enriched capability category.
For irreducible cooperative saturating networks, a robust reusable bit
exists for some positive process and read margins exactly when the Perron
root exceeds one. Rational positive and negative certificates make
finite-candidate topology synthesis effective. A four-probe test can also
certify a latch with unknown monotone transfer laws under explicit source
and noise promises. Constructed invariant regions support arbitrarily
delayed experiments that no fixed finite feedforward network can perform
with the same positive read margin. Manufacture precedes the future bit,
and no external bit-retaining controller is allowed.

The formation comparison distinguishes closure after joint composition
from separate formation, and separates semantic operations from native
witnesses. Reuse and transport require their own preservation conditions;
a categorical receptacle does not supply fabrication or empirical
validation. The result is a proved synthesis for stated native classes,
not a general solution to invention or a historical priority claim.
Algorithmic chemistry, process resource theories, dynamical circuit
theory, constructor theory, program synthesis, active learning, and
changing categorical interfaces supply direct antecedents.
\end{abstract}
\tableofcontents


% Opening module for the standalone manuscript.
\section{What changes when an interaction makes a tool}\label{sec:question}
An interaction can change more than the investigator's description of an
unchanged plant. It can assemble a catalyst, connect previously independent
modules, or install a device that supports a later experiment. The relevant
question is whether that construction supplies an executable and reusable
interface that was unavailable in the previous installed configuration.
Answering it requires a construction, a correctness argument, a resource
account, and a comparison. Merely declaring a new sort or admitting a new
semantic operation supplies none of these.

The target here is \emph{reachable installed capability growth}. Starting
with a stated stock and independently specified native interactions, a
finite manufacturing protocol produces an artifact. We derive a service
interface for the artifact and prove that subsequent uses implement that
interface. Future service inputs are not known during manufacture. The
result includes a task that can be carried out after installation and a
proof that the same service contract cannot carry it out beforehand.
All material inputs, consumed resources, retained states, and allowed
controllers belong to the comparison.

There are two constructive native classes in this note. Finite
mass-conserving reaction systems admit exhaustive synthesis and checking
of manufacture and reusable service cycles. A nonlinear circuit built
from individually fading modules supplies a robust autonomous bit. In
the second example, feedback changes the dynamics, creates separated
invariant information-bearing regions, and supports indefinitely delayed
experiments. The bit is not supplied as an acquisition token. Its write,
hold, and read operations are constructed and proved correct.

These constructions remain within their supplied substrate. They do not
escape the extensional closure of the complete factory runtime. Nor does
finite testing establish an eternal reliability guarantee without a
warranted model and disturbance bounds. Those limitations constrain the
meaning of the positive result; they do not turn fabrication into a
change of notation.

\subsection{Four different comparison questions}
Fix a native runtime, an initial physical configuration $q$, and a class
of admitted observations. We distinguish:
\begin{enumerate}[label=(\roman*)]
\item \emph{Global runtime expressivity}: which complete manufacture-and-use
policies can run from $q$?
\item \emph{Installed service capability}: which tasks can an already
installed configuration perform under a fixed service protocol?
\item \emph{Cost and availability}: how much time, stock, energy, storage,
or repeated preparation does a task require?
\item \emph{Warrant}: which claimed implementations are justified by the
available evidence and independently assumed source model?
\end{enumerate}
A strict answer to one question need not be a strict answer to another.
A manufactured tool may be newly available for service and reduce the cost
of many future uses, even though the factory program that manufactures
and uses it was executable from the beginning. Conversely, introducing a
new interface without a native realization may change a semantic language
without changing any physical capability.

The prior notes on native operations distinguished a supplied grammar,
its native realizations, semantic operations, and witnesses. The present
note adds an explicit state-changing manufacture step. It is therefore a
continuation of that distinction, not a replacement of it. Section~\ref{sec:project}
states exactly which previous constructions it uses and which stronger
claims remain open.

\section{The resource comparison and its closure boundary}\label{sec:boundary}
A configuration includes the artifact, stock, wiring, control state, and
all persistent memory that can influence service. A \emph{manufacturing
protocol} consists of allowed native actions, possibly chosen from earlier
observations, performed before the future service input arrives. A
\emph{service contract} $\mathcal S$ specifies its input and output ports,
permitted control actions, timing, retained state, disturbance model,
observations, and success condition. Manufacturing controls are available
during service either to both compared configurations or to neither.

Write $\operatorname{Cap}_{\mathcal S}(q)$ for the tasks implementable from
$q$ under that contract. This notation records actual realizations, not
merely syntactically well-typed programs. Reachable installed growth is a
native manufacturing execution
\[
                 m:q_0\longrightarrow q_1
\]
and a task $T$ for which
\[
 T\in\operatorname{Cap}_{\mathcal S}(q_1)
 \qquad\text{and}\qquad
 T\notin\operatorname{Cap}_{\mathcal S}(q_0).
\]
Set inclusion of \emph{all} capabilities requires a separate argument:
an irreversible manufacture may sacrifice other services. The existence
of a strictly gained task alone does not prove that
$\operatorname{Cap}_{\mathcal S}(q_0)$ is a subset of
$\operatorname{Cap}_{\mathcal S}(q_1)$.

\begin{theorem}[Compilation of a complete construction and use]\label{thm:closure}
Let $\mathcal P(q_0)$ contain every complete policy admitted by a fixed
native runtime from a fixed initial configuration $q_0$. Suppose an
extended interface is implemented by native representations and uniform
native programs for manufacture, initialization, each operation, and
retained-state handling. Suppose their observations and composition agree
with the extended interface semantics. Then every finite extended
experiment on old boundary ports is implemented by a member of
$\mathcal P(q_0)$. The statement extends to infinite or recursive
experiments only when the native translation preserves the corresponding
iteration and stopping semantics. A resource bound transfers only when
the translation preserves or explicitly bounds the relevant costs.
\end{theorem}
\begin{proof}
Translate each occurrence of a new operation to its native implementation
and each new value to its native representation. Translate initialization
and retained-state passage as part of the program, rather than treating
them as free constants. Induction over finite experiment formation shows
that the translated program has the same boundary observations. Prefix
that program with the native manufacturing implementation. Closure of the
runtime under the declared sequencing and branching rules puts the
complete program in $\mathcal P(q_0)$. Recursion or infinite interaction
cannot be justified by finite structural induction; the additional
preservation hypothesis gives precisely that step. Likewise an
extensional simulation does not imply equal cost, which is why the last
claim requires a separate cost premise.
\end{proof}

The theorem is an elementary compilation fact and an accounting rule.
Runtime code generation, feedback, self-modification, and learning do not
invalidate it when their actual implementations are already part of the
runtime. Established expressiveness theory supplies stronger and more
specific translation criteria; see Felleisen~\cite{Felleisen1991}.

\begin{proposition}[The matched service test]\label{prop:fairness}
A task-level gain after manufacture is consistent with
Theorem~\ref{thm:closure}. It is not a strict global expressivity gain if the
baseline is allowed the same manufacture followed by the same service.
If the manufacturer retains a copy of the future service input, or if
one side receives an uncharged controller that already implements the
task, that comparison does not establish a gain of the artifact alone.
\end{proposition}
\begin{proof}
The concatenated manufacture-and-service policy is an admitted baseline
policy under the global comparison, so it reproduces the after-manufacture
behavior. Under the installed comparison, the two initial configurations
differ while the permitted service actions remain fixed, and that
concatenated policy need not be an allowed service policy. A retained
input or uncharged implementing controller is part of the physical
configuration; omitting it changes the stated initial state and hence
the problem being compared.
\end{proof}

For the robust memory below, both configurations have the same write and
read ports, the same noise requirements, and the same prohibition on
control during the hold phase. The difference is a manufactured feedback
connection installed before the bit is supplied. The writing controller
is reset or removed before holding, and no external memory of the input
is available to the reader.

\subsection{What finite evidence can warrant}\label{sec:evidence}
Let $W$ be a supplied class of possible native worlds. An observation
protocol produces an evidence transcript $E(w)$ in world $w$. A proposed
artifact and its use protocol count as certified only if their advertised
contract holds in every world still compatible with the evidence and
source assumptions.

\begin{proposition}[Evidence fibre obstruction]\label{prop:fibre}
For each world $w$, let $I_w$ be the set of native implementations valid
for a requested contract. A sound evidence-only selection at transcript
$e$ must lie in
\[
                    \bigcap_{E(w)=e}I_w.
\]
If that intersection is empty, no procedure using only $e$ can return a
sound implementation. In particular, if two worlds with identical
evidence disagree on feasibility, evidence alone cannot decide
feasibility correctly in both.
\end{proposition}
\begin{proof}
The procedure has the same transcript in every world in the fibre and
therefore returns the same implementation. Soundness requires membership
in every corresponding $I_w$. The decision claim follows because the
same answer would be assigned to two different truth values. A randomized
procedure satisfying sure soundness has its output support in the same
intersection; randomization does not remove the obstruction.
\end{proof}

A finite state model, a polynomial differential law, a bounded parameter
interval, and a disturbance envelope are different forms of source
warrant. They may be justified independently, but they cannot be inferred
from arbitrary finitely many observations merely because a fitting model
exists. The construction theorems below state their model promises rather
than assuming that the desired artifact already has a valid interface.


% Source-ready module. Assumes amsmath, amssymb, theorem/proposition/corollary/example/remark environments.
% All results below concern a stipulated executable discrete reaction model, not an inferred chemical mechanism.
\section{Endogenous manufacture in a fixed executable substrate}

The substrate in this section is fixed throughout. The individual artifact, its fabrication history, and its cyclic service implementation are outputs of the construction, rather than extra primitives introduced after learning what service is wanted. Nevertheless, the elementary transition laws and the ability to execute them are premises. No theorem can infer an unobserved physical law solely from a name assigned to a manufactured object.

\subsection{Native states, elementary steps, and three comparison contracts}
Fix a finite set $\Sigma$ of species and a finite set $T$ of elementary instructions. An instruction $t$ is a pair of multisets $\alpha_t,\beta_t\in\mathbb N^{\Sigma}$. At marking $x\geq\alpha_t$, its specified native execution is
\[
 x\xrightarrow{t}x-\alpha_t+\beta_t.
\]
The specification includes the ability to select and complete this elementary instruction when enabled. It is not inferred from a stoichiometric equation alone. Spontaneous stochastic chemistry, adversarial reactions, inhibition, spatial separation, finite-precision detection, and wear require different runtime models or additional warranties.

A finite word of enabled instructions is a native program. Parallel independent programs are combined by multiset addition; identities and sequential composition give the usual process category of the reaction system. One can retain firing histories and interchange identifications, or pass to its thin reachability category. All existence and impossibility results below hold in either account: every finite compositional process has a sequential firing realization, and every such realization is a process.

Assume a supplied positive integral conserved grading
\begin{equation}\label{er:eq:mass}
 \mu:\Sigma\longrightarrow\mathbb N_{>0},\qquad
 \mu\mathbin\cdot\alpha_t=\mu\mathbin\cdot\beta_t\quad(t\in T).
\end{equation}
A grading can represent atom count or a formal conserved mass; it is not a claim that this one number is a complete physical model. For $m\in\mathbb N$, write
\[
 X_m=\{x\in\mathbb N^{\Sigma}:\mu\mathbin\cdot x=m\}.
\]
This is finite, with the safe bound $|X_m|\leq {m+|\Sigma|\choose |\Sigma|}$, because every such marking contains at most $m$ tokens. Let $G_m$ be the finite directed graph of enabled elementary steps on $X_m$.

A decidable predicate specifies which markings are stable resource stocks. It can, for example, exclude transient bound complexes and client input/output species. Fix an initial stable stock $q$ of mass $m$. A finite family of client service contracts is specified by input and output multisets
\[
 Q_i=(u_i,v_i),\qquad\mu\mathbin\cdot u_i=\mu\mathbin\cdot v_i=d_i.
\]
They are boundary specifications, not assumed implementation instructions. Define the ready, stock-preserving capability profile
\begin{equation}\label{er:eq:profile}
 \operatorname{Cap}(s)=\{i:s+u_i\longrightarrow^*s+v_i\}.
\end{equation}
The entire stock $s$ must be returned, and no undeclared material is supplied or discarded. A new service means an element of $\operatorname{Cap}(s)\setminus\operatorname{Cap}(q)$; strict repertoire growth means the stronger relation $\operatorname{Cap}(q)\subsetneq\operatorname{Cap}(s)$. Manufacturing need not preserve every old capability, so these are distinct tests. A service path may use intermediates and even fabrication steps if they can actually restore that stock.

There are three distinct comparisons.
\begin{enumerate}
\item A \emph{stationary service} requires the displayed resource stock to be returned. A strict increase compares $\operatorname{Cap}(q)$ with the profile of an actually manufactured stock $s$.
\item A \emph{first-use latency} comparison allows manufacture but charges its steps, and may allow the final internal stock to differ.
\item A \emph{complete experiment} includes manufacture, use, and all final internal material in one native process. It never escapes the closure of the fixed substrate.
\end{enumerate}
The first two are substantive operational comparisons. Neither implies an old-boundary behavior outside the third language.

\begin{theorem}[Exact endogenous installation and service synthesis]\label{er:thm:synthesis}
Under the finite executable model and conserved grading above, there is an effective algorithm which computes all stable stocks $s$ obtainable from $q$, together with their exact profiles $\operatorname{Cap}(s)$. For each reported capability it returns:
\begin{enumerate}
\item a native fabrication word $a:q\to s$;
\item a native cyclic service word $w_i:s+u_i\to s+v_i$;
\item whenever $i\notin\operatorname{Cap}(q)$, an independently checkable finite obstruction certificate for any stationary implementation at $q$.
\end{enumerate}
Every obtainable stable stock and every realizable declared stock-preserving service contract is detected. The algorithm returns a witness, not every execution history. In particular, it decides whether any obtainable stock has a strict profile gain, and produces its elementary fabrication and service witnesses when one exists. The acquired service need not occur as a primitive instruction or be specified in advance as the action of a named catalyst.

With nonnegative rational instruction costs, the same construction computes shortest fabrication and service costs, and an optimal installation for any specified finite number $n$ of identical calls within the class of plans which first install one fixed stock and return it after every isolated call. Each call receives only its specified client input, and its output is delivered before the next call; client material is not pooled across calls.
\end{theorem}
\begin{proof}
Run graph reachability from $q$ in $G_m$, retaining a predecessor for every visited marking. Intersect the reachable set with the decidable stable-stock predicate. Every retained marking $s$ has a reconstructible fabrication word, and every obtainable stable stock occurs because each elementary execution preserves mass.

For each such $s$ and contract $i$, run reachability in $G_{m+d_i}$ from $s+u_i$. Include $i$ in the profile precisely if $s+v_i$ is visited. Predecessors reconstruct the service word. Soundness is the enabledness check on every edge. For completeness, every finite native process sequentializes to a word of elementary firings and therefore appears as a path of the corresponding graph. This argument does not replace a firing sequence by a merely balanced stoichiometric vector: intermediate enabledness is checked at every step.

If the target is not reached, the reachable vertex set $Z$ is a certificate: it contains the source, excludes the target, and is closed under every elementary transition enabled at a member of $Z$. These three finite properties can be checked without trusting the search implementation. Induction on path length then proves the target unreachable. One can supply a smaller invariant or transition-closed set instead whenever available.

For costs, use finite-graph shortest-path algorithms; zero-cost cycles cause no difficulty for existence of a shortest finite witness because a repeated vertex can be removed without increasing cost. Write $A(q,s)$ for minimum fabrication cost and $L_i(s)$ for minimum cyclic service cost, with infinity for an unavailable path. For a fixed installed stock $s$, an $n$-call stationary plan costs at least $A(q,s)+nL_i(s)$ and achieves this bound by concatenating shortest witnesses. Minimize this finite list over obtainable stable $s$. The assertion is intentionally restricted to fixed-installation plans; a plan that changes its stock between calls belongs to a larger optimization problem.
\end{proof}

The algorithm is finite, not necessarily efficient. One direct implementation uses at most $|X_m|$ fabrication states and at most $|X_{m+d_i}|$ states for each service search, for each candidate stock. The positive grading is a checkable sufficient condition for this finite bound. Boundedness can also be established by other supplied invariants. No claim of comparable complexity or termination is made after removing every finite-state bound.

\begin{corollary}[Repeated use and material accounting]\label{er:cor:reuse}
If $a:q\to s$ and $w:s+u\to s+v$ are returned witnesses, then for every specified finite $n$ there is a native program
\[
 q+n u\longrightarrow^*s+n v.
\]
It consists of the fabrication word and $n$ copies of the service word. Its cost is the fabrication cost plus $n$ times the service cost. Alternatively, $n$ client inputs may be supplied and $n$ outputs delivered sequentially, with $s$ retained between calls.
\end{corollary}
\begin{proof}
Add the unused client material to each enabled firing of $a$ and of each service word. Monotonicity of multiset enabledness makes every lifted firing valid. Induct on the number of service invocations. The resulting total mass is $m+n d$, not $m$; a finite-state bound for the original fabrication stock does not supply infinitely many client substrates. In the sequential account, every external input and output is a separately tracked boundary transfer.
\end{proof}

\subsection{The exact categorical reset boundary}
Let $\mathcal P$ be any symmetric monoidal native process category. For a fixed resource object $R$, define a category with old client object labels and
\[
 \mathcal P_R(A,B)=\mathcal P(R\otimes A,R\otimes B).
\]
Its identity is $1_{R\otimes A}$ and its composition is that of $\mathcal P$. These are state-threaded operations; no copying or discard of $R$ is assumed. Independent resources give functors
\[
 \mathcal P_R\times\mathcal P_S\longrightarrow\mathcal P_{R\otimes S}
\]
by tensor and symmetry. There is in general no tensor operation sharing one linear resource between two supposedly simultaneous uses, and no canonical capability-transport functor induced by an arbitrary fabrication arrow $R\to S$.

\begin{proposition}[Well-typed reset prevents an existence gain]\label{er:prop:reset}
Suppose native arrows $a:R\to S$ and $b:S\to R$ both exist. If $\mathcal P_S(A,B)$ is nonempty, then $\mathcal P_R(A,B)$ is nonempty, witnessed by
\[
 (b\otimes1_B)\,f\,(a\otimes1_A).
\]
Consequently, a strict gain defined by an empty old service hom-set and a nonempty new one forces failure of a native return arrow $S\to R$. No inverse identities are needed for this existence assertion. If $a,b$ are inverse, conjugation gives an isomorphism of the complete state-threaded categories. Without those identities the sandwich construction need not be a functor or preserve the full behavior of an arbitrary value-sensitive service.
\end{proposition}
\begin{proof}
The displayed composite is well typed from $R\otimes A$ to $R\otimes B$. If the arrows are inverse, conjugation preserves identities and composition by cancellation, with inverse conjugation in the other direction. Without inverse identities a composite of two sandwiches inserts $ab$, so no general functor law follows.
\end{proof}

A reset which destroys the client output has type, for example, $S\otimes B\to R\otimes A$ or $S\otimes B\to R$, and is not the return map required by this proposition. A reset consuming fresh stock $F$ has type $S\otimes F\to R\otimes W$ with its actual waste $W$; it is not an uncharged arrow $S\to R$. In a full resource account, the extra stock and waste remain part of the endpoints. A valid reset preserving the output does eliminate the claimed unbounded stationary existence separation; any surviving advantage must then be one of cost, timing, reliability, or another explicitly restricted contract.

\subsection{An elementary coupled mechanism, with no primitive service macro}
Consider species $A,B,D,F,E,W,U,I,J,V$ and exactly these five elementary instructions:
\begin{align*}
 t_1:&\quad A+B\longrightarrow D,\\
 t_2:&\quad D+F\longrightarrow E+W,\\
 t_3:&\quad E+U\longrightarrow I,\\
 t_4:&\quad I\longrightarrow J,\\
 t_5:&\quad J\longrightarrow E+V.
\end{align*}
Here $D$ is an assembled intermediate, $F$ a one-use activation reagent, $W$ retained activation waste, and $I,J$ transient service complexes. These names aid interpretation but do not supply any additional operation. Assign weights
\[
 \mu(A)=\mu(B)=\mu(F)=\mu(W)=\mu(U)=\mu(V)=1,
 \quad\mu(D)=\mu(E)=2,
 \quad\mu(I)=\mu(J)=3.
\]
All five instructions conserve the grading. The model is a mathematically executable reaction mechanism, not a claim that particular chemical molecules realizing it have been built.

Let $q=A+B+F$, $s=E+W$, and $Q=(U,V)$. A fabrication word is $a=t_1t_2$. The service word $w=t_3t_4t_5$ is derived, rather than inserted as a primitive. Its actual stock trace is
\[
 E+W+U\xrightarrow{t_3}I+W
 \xrightarrow{t_4}J+W
 \xrightarrow{t_5}E+W+V.
\]
The reusable catalyst $E$ is physically absent at the beginning; the only way to obtain it from $q$ is to execute the assembly and activation interactions. Both the intermediate $D+F$ and the waste $W$ are accounted for.

\begin{theorem}[Strict endogenous stationary gain and exact setup cost]\label{er:thm:example}
For this fixed five-instruction system:
\begin{enumerate}
\item $U\to V$ is a cyclic service at $s$, but not at $q$, nor at the intermediate stock $D+F$.
\item The separate stocks $A+F$ and $B$ cannot produce $V$ from supplied $U$, even if their final resource stocks need not be restored. Their coupled stock $q$ can manufacture $s$ and then produce $V$.
\item For $n\geq1$, the shortest process $q+nU\to s+nV$ has exactly $2+3n$ firings. The shortest cyclic service process $s+nU\to s+nV$ has exactly $3n$ firings.
\item There is no reset $s\to q$. Hence this stationary separation is not defeated by an uncharged prepare--use--reset implementation. If valid reverse assembly and activation instructions are added, the stationary impossibility disappears, while charged setup and reset costs remain.
\end{enumerate}
\end{theorem}
\begin{proof}
The displayed service trace proves availability at $s$. For a hypothetical path $q+U\to q+V$, the numbers of $A,B,F$ must return to one. These species are never produced, while $t_1$ consumes $A,B$ and $t_2$ consumes $F$. Therefore neither $t_1$ nor $t_2$ can fire. Initially no $E,I,J$ is present, so $t_3,t_4,t_5$ cannot fire either. Thus $V$ cannot be produced. At stock $D+F$, preservation of $F$ similarly excludes $t_2$, and then excludes the entire service chain.

From $A+F+U$, no $B$ is present or producible; therefore $t_1$ cannot fire, and no $D,E,I,J,V$ can be produced. From $B+U$, the corresponding missing $A$ gives the same obstruction. Coupling supplies both assembly inputs; the word $t_1t_2t_3t_4t_5$ is then enabled.

For the cost statement, producing $n$ copies of $V$ requires exactly $n$ occurrences of $t_5$, since only $t_5$ produces $V$ and no instruction consumes it. Initially and finally $I,J$ are absent. Balance of $J$ forces $n$ occurrences of $t_4$; balance of $I$ forces $n$ occurrences of $t_3$. From $q$, the final absence of $A$ and presence of one $E$ force one $t_1$ and one $t_2$. From $s$, neither is enabled or needed. The constructed words meet the resulting lower bounds.

Finally, neither $A$ nor $B$ nor $F$ is produced by any instruction, so $s\to q$ is impossible. Adding $E+W\to D+F$ and $D\to A+B$ gives an output-preserving return from $s$ to $q$, so Proposition~\ref{er:prop:reset} constructs a stationary service at $q$ in seven firings.
\end{proof}

The total closed computation $q+U\to s+V$ already belongs to the original five-instruction language. The theorem does not deny this or pretend that $U\to V$ has escaped its global closure. What changes endogenously is the actual reusable internal stock and its exactly resource-preserving continuation repertoire. The $2+3n$ formula additionally shows exactly where the fabrication cost went.

\subsection{A tractable synthesis subclass and finite feedstock}
Let $K$ be a finite set of catalyst labels and $\mathcal H$ a finite collection of Horn rules $D\Rightarrow j$, with $D\subseteq K$. For each rule supply the native fabrication instruction
\begin{equation}\label{er:eq:horn}
 F+\sum_{i\in D}C_i\longrightarrow C_j+\sum_{i\in D}C_i.
\end{equation}
There is one common feedstock species $F$, consumed once per firing. A premise catalyst is retained, not duplicated; the displayed identical input and output occurrences mean that the same catalytic stock is restored by the elementary instruction. A physical implementation must warrant this behavior. Every premise occurs at most once; rules requiring multiplicity need a multiset closure instead. Give every $F,C_i$ weight one, so fabrication conserves total mass.

For an initial support $S\subseteq K$ and integer $b\geq0$, put $M(S,b)=\sum_{i\in S}C_i+bF$. Let $\operatorname{cl}(S)$ be the least Horn-closed subset of $K$ containing $S$.

\begin{theorem}[Exact finite-stock realization of Horn closure]\label{er:thm:horn}
Every catalyst label producible from $M(S,b)$ lies in $\operatorname{cl}(S)$. If $b\geq|\operatorname{cl}(S)\setminus S|$, a native saturation algorithm constructs exactly
\[
 M\bigl(\operatorname{cl}(S),b-|\operatorname{cl}(S)\setminus S|\bigr)
\]
using exactly $|\operatorname{cl}(S)\setminus S|$ fabrication firings. This is the minimum number of firings needed to obtain the full closure support, and that feedstock bound is necessary. There is no assumption of reusable or infinite fuel. With less stock, Horn closure remains an upper bound but need not be physically achievable.
\end{theorem}
\begin{proof}
Induction on firings proves the first assertion: the only newly produced label is the head of a rule all of whose premises are already present. For saturation, repeatedly choose an enabled rule whose head is not yet present. Such a rule exists until the current support is closed; the least-closure property then identifies termination support with $\operatorname{cl}(S)$. Each firing preserves every previously present catalyst, adds one absent label, and consumes exactly one $F$. The assumed initial stock therefore suffices at every stage. Each missing label requires a firing producing it, and each firing produces at most one new label. Thus at least the stated number of firings and feedstock tokens is required.
\end{proof}

To connect this closure to reusable services without treating a service macro as primitive, one may use, for each $i$, the uniform elementary mechanism
\[
 C_i+U_i\to I_i,\qquad I_i\to J_i,\qquad J_i\to C_i+V_i.
\]
Give $U_i,V_i$ weight one and $I_i,J_i$ weight two. There are no other service instructions. These are fixed substrate rules, not resources initially available. For stock $M(S,b)$ the exact stationary service profile is $S$: preserving the number of $F$ excludes every fabrication firing; if $i\notin S$, neither $C_i,I_i,J_i$ can then appear, whereas if $i\in S$ the three-step cyclic implementation is enabled. The Horn theorem thus compiles an abstract closure to actual installed capability growth in this explicit native class. It does not establish that an arbitrary real substance has the specified elementary reaction mechanism.

\subsection{Finite evidential bridge with charged sacrificial specimens}
The algorithm above takes a warranted native model as input. Finite observations alone do not warrant an unrestricted reaction model. The following limited bridge states exactly what additional experimental assumptions suffice; it does not silently reset a manufactured artifact for free.

\begin{proposition}[Finite identification with fully observed deterministic copies]\label{er:prop:identify}
Suppose a substrate has a finite, fully observed state space of size at most $N$, a known set of $k$ elementary action labels, deterministic action effects, and a non-destructive blocked response for a disabled action. Assume identical fresh specimens can be prepared in one specified initial state, and the action law is stationary across specimens. A complete reachable transition table, together with replay words for its states, can be learned using at most $kN$ sacrificial specimens and at most $kN^2$ elementary action attempts. A further fresh specimen can execute a chosen returned plan. If there are several initial test configurations, apply and charge this bound separately to each configuration.
\end{proposition}
\begin{proof}
Maintain a breadth-first list of observed states and one shortest replay word from the initial state to each. For every discovered state and every action label, take a fresh specimen, replay that state's word, and attempt the label. Record either its observed successor or its blocked status, adding a new state when necessary. At most $N$ distinct states occur, so there are at most $kN$ such tests. Each shortest replay has length at most $N-1$, and the tested action adds at most one attempt, giving the bound. When the queue is exhausted, every action at every reached state is known. Induction on execution length proves the recorded table contains every reachable state and transition. Determinism and cross-specimen stationarity justify replay; exact full observation justifies identifying repeat states.
\end{proof}

These are substantial model and access assumptions. Finite tests can reveal violations but do not certify their unrestricted validity. The consumed specimens, any client substrate used in them, and their final wastes must be included in the experimental resource ledger. Preparing a new specimen is not a return arrow from the old specimen's final state to its initial state. If a genuine in-place, output-preserving native reset is available instead, Proposition~\ref{er:prop:reset} applies and the stationary existence separation must be reconsidered. With destructive single-shot access, hidden state, nonstationarity, or no finite bound, this identification theorem does not apply. The mathematical implication from a complete observed transition table to the reported finite plans is exact; transfer from the specimens to the real device remains the stated empirical model warrant.

\subsection{Discovering a whole bounded capability doctrine, rather than one named service}
A finite target list is not essential. Choose a finite set of admitted boundary species $\Gamma\subseteq\Sigma$ and a finite client mass bound $L$. Let $\mathcal B_L$ be the finite set of multisets supported on $\Gamma$ of mass at most $L$. These are the fixed old boundary sorts; no desired conversion between them is assumed. For an obtainable resource stock $s$, define
\[
 d_s(u,v)=\inf\{\operatorname{cost}(w):w:s+u\to s+v\},
 \qquad u,v\in\mathcal B_L,
\]
with infinity for no implementation. Conserved mass makes this infinity when $u,v$ have different masses. Nonnegative rational costs and finite-state boundedness make every finite infimum a witnessed minimum.

\begin{corollary}[Exact generated capability categories]\label{er:cor:doctrine}
The synthesis algorithm computes the full matrix $d_s$ for every obtainable stable stock $s$, together with a native shortest witness for every finite entry and a finite obstruction certificate for every infinite entry. These matrices satisfy
\[
 d_s(u,u)=0,\qquad d_s(u,w)\leq d_s(u,v)+d_s(v,w).
\]
Thus each is an extended asymmetric Lawvere pseudometric, equivalently a category enriched in the usual reversed-order additive cost quantale. Forgetting finite costs gives a preorder of native implementability on the unchanged boundary objects. The calculation discovers all bounded services rather than selecting only an already named catalyst action.
\end{corollary}
\begin{proof}
Enumerate the finite set $\mathcal B_L$ and use every equal-mass pair as a service contract in Theorem~\ref{er:thm:synthesis}. Identities are empty firing words. Composing shortest words through the exactly restored stock $s$ gives the triangle inequality. If either operand cost is infinite, the inequality is immediate. The same composition proves transitivity of finite reachability, and empty words prove reflexivity. Each returned finite entry is still accompanied by a native firing history; passing to the enriched matrix is an explicit abstraction, not a reconstruction of all proof-relevant histories.
\end{proof}

This bounded doctrine can contain newly finite entries after manufacture, as well as decreased finite costs. Pointwise $d_s\leq d_q$ defines an across-the-board cost improvement and makes the identity object map an enriched functor from the old doctrine to the new one. This inequality need not hold after arbitrary manufacture: some old services may become slower or unavailable. The algorithm detects gains, losses, and tradeoffs rather than assuming monotonicity. No arbitrary service oracle is placed into its definition: its entries are computed from elementary executions. The boundary grammar, the material bound, and the physical laws remain supplied. The object set is generally not closed under tensor beyond $L$; only combinations respecting the declared bound may be treated as bounded parallel tasks. Removing the bound requires a separately justified analysis.


% Standalone section fragment. Uses ordinary amsmath/amsthm theorem environments.
% This is an elementary closure/provenance synthesis, not a historical novelty claim.
\section{Formation, coupling, and certified interaction}
\label{syn-sec:formation}

Two orders must be distinguished. A stock can be ordered by the resources
available to a constructor. An installed device is instead a particular
physical state or arrangement, and changing it can remove as well as add
services. The first order often supports a closure operator. The second
does not do so merely because the first does. The results in this section
concern the first order, and transfer to installed capabilities only through
the independently proved manufacture and service theorems of an instance.

\subsection{An exact coupling law from primitive formation rules}

Let $K$ be a finite set of resource species and let $\mathcal R$ be a finite
set of implications $P\Rightarrow c$, with $P\subseteq K$ and $c\in K$.
Let $C(A)$ be the least superset of $A\subseteq K$ closed under these rules.
Thus $C$ is extensive, monotone, and idempotent. This is ordinary finite Horn
closure. A physical interpretation requires its own theorem: in the
catalytic instance, the premises must remain available after a conversion,
enough feedstock must be supplied, and every inferred species must be
manufactured by the stated elementary reactions. Horn syntax alone gives
none of those conclusions. Species presence does not count simultaneous
copies or parallel throughput.

Write $\mathcal L=\{D\subseteq K:C(D)=D\}$. In this lattice, meets are
intersections and joins are $D\vee E=C(D\cup E)$.

\begin{theorem}[Formation--coupling comparison and exact obstruction]
\label{syn-thm:coupling}
For initial stocks $A,B\subseteq K$, put $U=C(A)\cup C(B)$.
Then
\[
 C(A\cup B)=C(C(A)\cup C(B))=C(A)\vee C(B).
\]
The canonical comparison from the separately formed resource sets to the
jointly formed resource set is the inclusion
\[
 U\ \subseteq\ C(A\cup B).
\]
It is an equality if and only if every primitive rule satisfies
\[
             P\subseteq U\quad\Longrightarrow\quad c\in U.
\tag{*}
\]
If equality fails, a primitive rule $P\Rightarrow c$ witnesses the failure
with $P\subseteq U$, $c\notin U$, and
\[
 P\cap(C(A)\setminus C(B))\ne\varnothing,
 \qquad
 P\cap(C(B)\setminus C(A))\ne\varnothing.
\]
Its conclusion is therefore a certified first interaction-dependent
resource, rather than a resource manufactured by either component alone.
\end{theorem}
\begin{proof}
Monotonicity gives $C(A)\cup C(B)\subseteq C(A\cup B)$, so closing again
gives one inclusion. Since $A\cup B\subseteq U$, the opposite inclusion
also holds. The join formula follows. Equality holds exactly when $U$ is
closed, which is equivalent to the primitive check $(*)$.

If equality fails, $U$ is not closed. Choose a violated primitive rule.
If its premise were contained in $C(A)$, closure of $C(A)$ would put its
conclusion there; likewise for $C(B)$. The premise is contained in their
union but in neither separately, proving the two displayed intersections.
\end{proof}

The criterion is finite and input-level: compute two closures and inspect
each primitive implication. It is not a premise asserting that all joint
resources are already separately realizable. It gives both an exact
positive test and an explicit negative certificate. Repeated coupling
obeys
\[
 C(C(A\cup B)\cup D)=C(A\cup B\cup D)
                    =C(A\cup C(B\cup D)).
\]
This is associativity of formation with complete retained inventories;
it is not a claim that irreversible physical execution orders, preparation
costs, or installed devices are interchangeable.

The same theorem allows shared factory resources $F$: replace $A,B$ by
$F\cup A,F\cup B$. No copy of $F$ is inferred. The Boolean statement
concerns sequential availability with persistent prerequisites. Any
simultaneous use of one factory resource needs its own capacity account.

\subsection{Minimal supports retain what a bare service list forgets}

For $c\in K$ define its family of minimal supports
\[
 \mathcal M_c=\min_{\subseteq}\{S\subseteq K:c\in C(S)\}.
\]
This family includes $\{c\}$ unless a strictly smaller support suffices.
Only supports contained in the actual initial stock are usable; the
singleton does not authorize assuming the resource one seeks to make.

\begin{theorem}[Exact support and interaction certificates]
\label{syn-thm:support}
For every $A\subseteq K$,
\[
            c\in C(A)\quad\Longleftrightarrow\quad
            (\exists S\in\mathcal M_c)\ S\subseteq A.
\]
The families $\mathcal M_c$ can be constructed from the primitive rules by
the following terminating procedure. Start with the support $\{c\}$ at
each $c$. For each rule $P\Rightarrow c$, enumerate every tuple of currently retained
supports $S_p$, one for each $p\in P$, and insert $\bigcup_{p\in P}S_p$
at $c$. Discard supports strictly containing another retained support.
For an empty premise insert the empty support. Repeat complete passes over
all rules and support tuples until a pass makes no change.
Every retained support can carry a finite derivation certificate.

In particular, $c$ belongs to
$C(A\cup B)\setminus(C(A)\cup C(B))$ if and only if there is a support
in $\mathcal M_c$ contained in $A\cup B$, and every such support meets
both $A\setminus B$ and $B\setminus A$.
\end{theorem}
\begin{proof}
The first assertion follows by taking an inclusion-minimal subset of a
finite generating stock. Every support inserted by the algorithm is sound:
execute the certified derivations of the premises and then the rule.
Discarding a larger support cannot prevent a later rule application,
because replacing it by a smaller support only shrinks the resulting
union. More formally, at every step take the upward closure of each
retained support family. These upward-closed sets only increase, and
there are finitely many subsets of $K$, so the process stabilizes.

At stability, every finite rule derivation is represented: induction on
its last rule shows that the union of its premise supports is covered by
a retained support. Conversely all retained supports have derivations.
The minimal members are exactly $\mathcal M_c$. Finally $c\notin C(A)$
means that no usable minimal support is contained in $A$, and similarly
for $B$. A subset of $A\cup B$ lies in neither $A$ nor $B$ exactly when
it meets both displayed exclusive parts.
\end{proof}

This provenance is stronger than a Boolean claim that a capability exists.
It distinguishes alternative routes, shared requirements, and unavoidable
cross-component dependence. The antichains can be exponentially large;
the theorem promises decidability and certificates, not polynomial-size
provenance. Nor does a smallest support necessarily minimize reaction
count, elapsed time, fuel, or simultaneous storage.

For example, take rules $\{a,b\}\Rightarrow c$ and
$\{c,d\}\Rightarrow m$, with $A=\{a,d\}$ and $B=\{b\}$.
Separate formation yields no $c$ or $m$, whereas joint formation makes
both. The first cross-rule produces $c$, and the usable minimal support
for $m$ is $\{a,b,d\}$. By contrast, with rules
$a\Rightarrow c$ and $b\Rightarrow d$, the separate union is already
closed. Even a syntactically mixed rule need not produce a genuine gain:
if $a\Rightarrow c$ is also present, the rule
$\{a,b\}\Rightarrow c$ is interaction-redundant.

\subsection{When a change of doctrine preserves the generated gain}

Let a second finite rule presentation have species $K'$, rules
$\mathcal R'$, and closure $C'$. Let $f:K\hookrightarrow K'$ be an
injective interface interpretation. The target may have auxiliary
species and different elementary formation rules. Source-rule preservation
means the finitely checkable condition
\[
   (P\Rightarrow c)\in\mathcal R
       \quad\Longrightarrow\quad f(c)\in C'(f[P]).
\tag{P}
\]
For physical use this condition must additionally carry a native target
recipe for the target entailment. It is not enough to map species labels.

\begin{theorem}[Recognizable conservative formation transport]
\label{syn-thm:transport}
Condition $(P)$ gives $f[C(A)]\subseteq C'(f[A])$ for all $A$.
Under $(P)$, the following are equivalent:
\begin{enumerate}
\item For every $A\subseteq K$,
$f^{-1}[C'(f[A])]=C(A)$.
\item For every $c\in K$ and every target minimal support
$S\in\mathcal M'_{f(c)}$ contained in $f[K]$,
\[
                 c\in C(f^{-1}[S]).
\tag{R}
\]
\end{enumerate}
Thus preservation and reflection have finite certificates computed from
primitive rules and minimal supports, even when the target uses auxiliary
intermediates. When these conditions hold,
\[
 c\in C(A\cup B)\setminus(C(A)\cup C(B))
 \Longleftrightarrow
 f(c)\in C'(f[A\cup B])
            \setminus\bigl(C'(f[A])\cup C'(f[B])\bigr).
\]
Successive conservative interpretations compose conservatively.
\end{theorem}
\begin{proof}
Induction on a source derivation, replacing each source rule by its target
derivation, proves preservation. Reflection implies $(R)$ by applying it
to $f^{-1}[S]$, whose image is $S$.

Conversely, if $f(c)\in C'(f[A])$, choose a target minimal support
$S\subseteq f[A]$. Condition $(R)$ gives
$c\in C(f^{-1}[S])\subseteq C(A)$. This proves reflection. The gap
equivalence applies preservation and reflection to each of $A$, $B$,
and $A\cup B$. For composable conservative maps $f,g$, calculate
\[
 (gf)^{-1}C''(gf[A])
       = f^{-1}\bigl(g^{-1}C''(g[f[A]])\bigr)
       = f^{-1}C'(f[A])=C(A).
\]
Primitive preservation and its compiled witnesses compose by substitution.
\end{proof}

The auxiliary-species allowance matters: a target implementation may
realize a source formation rule through several intermediates without
having the source rule as one primitive. Conversely preservation alone
does not preserve a strict gain. In the preceding example, add a target
rule $a\Rightarrow m$. All source rules still compile, but $m$ is now
available to $A$ alone. The target support $\{a\}$ fails $(R)$ and
pinpoints the lost baseline obstruction.

Theorem~\ref{syn-thm:transport} is a finite construction-level counterpart
of the earlier native-reflection theorem across semantic doctrines.
That earlier theorem uses independent finite recognizers and source-test
coverage; here independently computed support antichains do the
recognition work. Neither theorem turns arbitrary semantic maps into
physical implementations, and neither recovers discarded execution
histories from extensional equality.

\subsection{What the common comparison establishes}

Finite Horn closure, implicational bases, closure-system lattices and
minimal-support provenance are established mathematical machinery.
The value of these results here is their exact connection to independently
proved native manufacture, stationary service semantics, and an unchanged
construction stock. Historical novelty is not claimed for the closure
theorems.

The continuous feedback latch is not silently another Boolean Horn model.
Its manufactured output is a changed installed vector field. The new
memory guarantee follows from dynamical stability, disturbance bounds,
write/read protocols, and a baseline impossibility proof, not from
closure of species labels. Both instances distinguish constructor access
from installed services, supply actual intermediate executions, and
exhibit a strict service comparison under the same postconstruction
tests. Their shared mechanism is native construction changing the
conditions of later use. The specific formation laws, monotonicity,
evidence, and robustness warrants remain doctrine-dependent.

There is no contradiction with whole-runtime conservativity. When the
constructor and subsequent uses are all included, both are executions of
the originally admitted runtime. A strict installed-interface result says
that the initial device state cannot supply a service satisfying the
specified deployment contract, whereas the constructed state can. It
does not say that the entire runtime computes a function outside its own
closure. Reset, external controllers, retained training data, feedstock,
copying, and online costs must be accounted for at the phase where they
are allowed.


\section{Feedback manufactures a robust memory interface}
\label{sec:feedback}
The reaction construction computes exact finite service paths. We now
construct a continuous-state interface whose reuse guarantee rests on
robust invariance rather than exact microscopic reset. This supplies a
qualitative new service from individually fading components and a
whole-baseline impossibility theorem.

\subsection{Native components and the service contract}
The native stock includes the leaky cells, memoryless gain and summing
stages, externally commanded selectors, unused connector terminals, a
wiring actuator, finite pulse timing, and bounded-error nondestructive
state reads. Uninterrupted power during a claimed operating interval is
an explicit common boundary resource for both the baseline and the
constructed device. No energy-optimality or unlimited physical-lifetime
claim is made. These primitives do not include a persistent bit.

Put
\[
 \sat(z)=\max(-1,\min(1,z)),
\]
acting coordinatewise on vectors. A primitive leaky cell with input signal $u$ has state $x\in\R$, output $x$, and dynamics
\begin{equation}\label{fb:eq:cell}
 \dot x=-x+\sat(u)+\eta.
\end{equation}
The external disturbance $\eta$ is measurable and satisfies $|\eta(t)|\leq\varepsilon$ almost everywhere. A gain-$g$ input stage is the memoryless operation $u\mapsto gu$. Under identical input and disturbance signals, two states of a single cell satisfy
\[
 |x(t)-x'(t)|=e^{-t}|x(0)-x'(0)|.
\]
Thus the cell itself forgets its history exponentially.

A finite coupled device is specified by a nonnegative real matrix $A$ and the public command $q\in\{\mathsf H,\mathsf W_+,\mathsf W_-\}$. Its dynamics are
\begin{equation}\label{fb:eq:network}
 \dot x=
 \begin{cases}
  -x+\sat(Ax)+\eta,&q=\mathsf H,\\
  -x+\boldsymbol 1+\eta,&q=\mathsf W_+,\\
  -x-\boldsymbol 1+\eta,&q=\mathsf W_-.
 \end{cases}
\end{equation}
A write command is a bounded-duration pulse, not an extra state variable holding the chosen bit forever. After either write pulse, the input is the same public hold command $\mathsf H$. The cells are the only persistent physical state in the modeled device. Nonnegative summing, gains and input selection are memoryless native operations. Their required transfer laws must be warranted independently of the memory theorem.

For fixed $\varepsilon\geq0$ let $M=1+\varepsilon$ and $E=[-M,M]^n$. Every command leaves $E$ invariant: at $x_i=M$, its derivative is at most $-M+1+\varepsilon=0$, and at $x_i=-M$ it is at least $M-1-\varepsilon=0$. All vector fields are globally Lipschitz in the state, and bounded measurable disturbances and piecewise constant commands have unique global absolutely continuous solutions. Disturbance and command changes do not reset the state. The supplied
initial stock has an arbitrary state in $E$; no bit value or invariant
memory region is initially promised.

\begin{definition}[Robust indefinitely reusable bit]\label{fb:def:memory}
A deployed device realizes such a bit if there are nonempty disjoint state regions $B_-,B_+\subset E$, a finite write time $T$, and a read margin $\nu>0$ with these properties:
\begin{enumerate}
\item For either sign $b$, the write pulse $\mathsf W_b$ of duration $T$ sends every initial state of $E$ into $B_b$, for every admitted disturbance.
\item Under the common command $\mathsf H$, every state of $B_b$ remains in $B_b$ for every finite hold duration, with the same nonzero disturbance bound.
\item A current-state read, with unknown additive observation error of sup norm at most $\nu$, determines $b$ correctly throughout $B_b$. Reading does not change the state.
\end{enumerate}
The read operation may be any decoder, but it cannot use an externally retained copy of the write bit or a previously stored observation. Hold durations have no predeclared upper bound. The number of cells and the error margins are fixed once the device is built.
\end{definition}

Positive read error is not cosmetic. With exact real-valued sensing, even $\dot x=-x$ can encode a sign forever in an arbitrarily small residue. Such an exact-sign convention is deliberately not the operational capability being claimed. A robust read must succeed on every state in a fixed-radius neighborhood, before any read-side amplification or computation.

\subsection{A sharp formation criterion for a parameterized class}
\begin{theorem}[Spectral threshold for robust memory]\label{fb:thm:threshold}
Let $A\geq0$ be an irreducible $n\times n$ matrix, let $\lambda=\rho(A)$, and use the native dynamics \eqref{fb:eq:network}.
\begin{enumerate}
\item If $\lambda>1$, explicit finite write controls construct a robust indefinitely reusable bit for some fixed positive process disturbance bound and read margin.
\item If $\lambda\leq1$, no such bit can be realized for any positive process disturbance bound and positive read margin, even if its proposed write preparation and decoder are arbitrary.
\end{enumerate}
Consequently the existence of this robust memory capability in the stated cooperative saturating class is equivalent to $\rho(A)>1$.
\end{theorem}
\begin{proof}
For sufficiency, choose a positive Perron vector $v$ with $Av=\lambda v$ and $\max_i v_i=1$. Set
\[
 a=\frac{1}{2\lambda},\qquad
 m=(\lambda-1)a\min_i v_i>0.
\]
Choose
\begin{equation}\label{fb:eq:eps}
 0<\varepsilon<\min\{m,1-a\},\quad M=1+\varepsilon,
\end{equation}
and define the two boxes
\[
 B_+=\prod_i[av_i,M],\qquad
 B_-=-B_+.
\]
If $x\in B_+$ and $x_i=av_i$, monotonicity gives
\[
 (Ax)_i\geq\lambda av_i,\qquad
 \sat((Ax)_i)\geq\lambda av_i,
\]
because $\lambda av_i\leq1/2$. Thus at that lower face
\[
 \dot x_i\geq(\lambda-1)av_i-\varepsilon\geq m-\varepsilon>0.
\]
At the upper face $x_i=M$, the derivative is at most zero as above. The standard first-exit argument, or the tangent-cone invariance criterion for a rectangle, proves that $B_+$ is invariant under every admitted measurable disturbance. Odd symmetry gives invariance of $B_-$. The inequality is strict on the inner faces, so this is not marginal stability.

During a positive write, variation of constants gives, for every $x_i(0)\geq-M$,
\[
 x_i(t)\geq1-\varepsilon+
   (x_i(0)-(1-\varepsilon))e^{-t}
 \geq1-\varepsilon-2e^{-t}.
\]
Since $av_i\leq a$, every duration
\begin{equation}\label{fb:eq:write-time}
 T\geq\log\frac{2}{1-\varepsilon-a}
\end{equation}
puts the whole state in $B_+$. The denominator is positive by \eqref{fb:eq:eps}; the negative write is symmetric. Reading the sign of any one coordinate $j$, with error $\nu<av_j$, decodes the bit correctly for all subsequent holds. For a common margin valid for every coordinate choose $0<\nu<a\min_i v_i$. This constructs all the data in Definition~\ref{fb:def:memory}.

For necessity, fix any $\varepsilon>0$, let $M=1+\varepsilon$, and put $\kappa=\varepsilon/M$. Choose the globally bounded disturbance law $\eta_i=-\kappa\max(-M,\min(M,x_i))$. Because $E$ is invariant, its realized path is
\[
 \eta(t)=-\kappa x(t).
\]
This is admissible globally, and along each run $|\eta_i(t)|\leq\varepsilon$. Equivalently, for each resulting solution its disturbance path is one of the bounded measurable paths over which the memory guarantee quantifies. Under this disturbance the hold system becomes
\[
 \dot x=-(1+\kappa)x+\sat(Ax).
\]
Since $A$ is nonnegative and $|\sat(z)|\leq|z|$, the coordinatewise upper right Dini derivative obeys
\[
 D^+|x|\leq-(1+\kappa)|x|+A|x|.
\]
For the positive Perron vector $v$, define the weighted maximum norm
$V(x)=\max_i|x_i|/v_i$. At every active maximizing coordinate,
$(A|x|)_i\leq\lambda V(x)v_i$, so
\[
 D^+V(x(t))\leq(\lambda-1-\kappa)V(x(t))
 \leq-\kappa V(x(t)).
\]
Therefore every initial state converges to zero exponentially under this admissible hold disturbance.

Take one state resulting from a putative write of each bit and apply the indicated admissible disturbance to each. The two resulting states become arbitrarily close. For any proposed positive read margin $\nu$, at a sufficiently large common read time their sup-norm distance is less than $2\nu$. Their observation-error cubes then intersect. The same noisy observation would have to be decoded both ways, which is impossible. The two disturbance histories need not coincide: robustness requires correctness for every admissible disturbance in either run. This contradiction applies to arbitrary preparation and arbitrary current-state decoders, not only to the particular boxes used in the positive construction.
\end{proof}

\begin{remark}[What the criterion does and does not characterize]
The equivalence is for the existence of \emph{some} positive process and read margins in this irreducible cooperative class. It does not say $\rho(A)>1$ suffices for any prescribed noise level, nor that every possible feedback architecture is governed by this Perron test. Irreducibility is used to obtain a strictly positive Perron vector. A reducible matrix can be treated through its strongly connected components, but that extension is not needed for the theorem as stated. At the threshold $\rho(A)=1$, noiseless continua of equilibria can occur; the necessity proof explains why they fail the demanded positive-noise guarantee.
\end{remark}

The constructive proof supplies an elementary invariant-box certificate. Exact knowledge of eigenvectors is not essential to execution: any positive rational vector $b$ with $\sat(Ab)-b$ componentwise bounded below by a certified positive number yields the same proof, with $B_+=\prod_i[b_i,M]$. Such strict inequalities can be verified using certified coefficient intervals. The experimenter need not measure an infinitely precise eigenvector or write it into the device. Its role is to warrant a finite safety margin.

\subsection{An effective classifier and a native wiring synthesis algorithm}
\label{sec:wiring-synthesis}
The spectral theorem gives more than a sufficient example. For rational
primitive data it yields a terminating certificate procedure. Thus a
constructor need not be handed a successful memory architecture or an
infinitely precise eigenvector.

\begin{theorem}[Finite rational certificates for topology synthesis]
\label{thm:rational-synthesis}
Let $A$ be an irreducible nonnegative rational matrix. Exactly one of the
following finite rational certificates exists:
\begin{align}
 \text{positive:}\quad&b\in\mathbb Q_{>0}^{n},
                \qquad b<Ab<\boldsymbol1;\label{eq:positive-cert}\\
 \text{negative:}\quad&w\in\mathbb Q_{>0}^{n},
                \qquad Aw\leq w.\label{eq:negative-cert}
\end{align}
All vector inequalities are coordinatewise. The positive case is
equivalent to $\rho(A)>1$, and the negative case to $\rho(A)\leq1$.
Dovetailing enumeration of rational vectors and these rational inequality
checks therefore terminates with a certificate of the correct case.

From a positive certificate one can compute rational $\varepsilon,\nu>0$
and an integer write duration $T$ realizing the robust memory contract.
Consequently, given a finite list of admissible rational irreducible
wiring matrices, together with a finite native link program constructing
each from the available modules and connectors, one can decide which
candidates implement some positive-margin robust bit, return their native
construction and service programs, or certify that none in the list does.
\end{theorem}
\begin{proof}
Let $p>0$ be a left Perron vector. If \eqref{eq:positive-cert} holds, then
\[
 \rho(A)p^{\mathsf T}b=p^{\mathsf T}Ab>p^{\mathsf T}b,
\]
so $\rho(A)>1$. If \eqref{eq:negative-cert} holds, the same calculation
with $w$ and a weak inequality gives $\rho(A)\leq1$. Hence the two cases
cannot both occur.

For $\rho(A)>1$, scale a positive right Perron vector $v$ so that
$Av<\boldsymbol1$. Then $Av>v>0$. All these inequalities are strict,
so a sufficiently close rational vector $b$ retains them. For
$\rho(A)<1$, take
\[
          w=(I-A)^{-1}\boldsymbol1
           =\sum_{k\geq0}A^k\boldsymbol1.
\]
The series converges, $w\geq\boldsymbol1>0$, and $Aw=w-\boldsymbol1<w$.
Because $I-A$ has rational entries and is invertible, $w$ is rational.
For $\rho(A)=1$, the rational kernel of $A-I$ has a rational basis.
Its real span contains a strictly positive Perron vector. Approximate
its real coefficients in that basis by rationals closely enough to
preserve strict positivity. The resulting rational $w>0$ satisfies
$Aw=w$. This proves existence in all cases. A fair enumeration eventually
encounters the existing certificate, and every check uses finitely many
rational operations.

Given $b$ in the positive case, let
\[
 d=\min_i((Ab)_i-b_i)>0,\qquad \beta=\max_i b_i<1.
\]
Choose a rational $0<\varepsilon<\min\{d,1-\beta\}$ and
$0<\nu<\min_i b_i$. The boxes
$B_+=\prod_i[b_i,1+\varepsilon]$ and $B_-=-B_+$ satisfy the inward
inequalities from Theorem~\ref{fb:thm:threshold}. Choose an integer $T$
such that
\[
                 \frac{2}{1+T}\leq1-\varepsilon-\beta.
\]
Since $e^T\geq1+T$, the lower write bound
$1-\varepsilon-2e^{-T}$ is at least $\beta$, so the entire box is
reached. The chosen $\nu$ is a valid coordinate read margin.

Run the classifier for every member of the finite candidate list. For
positive candidates concatenate the supplied elementary link program with
the computed write, hold, and read implementation. For negative
candidates, \eqref{eq:negative-cert} is a finite checkable obstruction:
the weighted-norm damping argument in the spectral theorem drives every
prepared state to zero under an allowed disturbance. A finite list is
exhausted by finitely many terminating classifications.
\end{proof}

The search space is supplied independently: a bounded inventory of
modules, a finite set of legal connector placements, and rational gain
values determine its candidates and link programs. The algorithm constructs
a successful architecture from that space; it does not assume one exists.
It decides neither unbounded circuit synthesis nor unrestricted physical
realizability. Its cost may be large, and no polynomial complexity bound
is asserted. Irreducibility is part of the stated classifier; the
feedforward impossibility result addresses a different, reducible class.

\begin{proposition}[Interval certificates and a sound acceptance test]
\label{prop:interval-warrant}
Suppose every deployed hold matrix is measurable, locally essentially
bounded, nonnegative, and obeys
$A(t)\geq A_-$ entrywise, where $A_-$ is a certified rational matrix.
Let $b\in\mathbb Q_{>0}^n$ satisfy $b<\boldsymbol1$ and
\[
       \sat(A_-b)-b\geq d\boldsymbol1\qquad(d>0).
\]
Assume the write field remains the commanded field in
\eqref{fb:eq:network}. For any
$0<\varepsilon<\min\{d,1-\max_i b_i\}$, the same boxes, read margin,
and finite write-time construction give a robust memory uniformly over
all these time-varying hold matrices and all admitted disturbances.
If the available coefficient intervals admit no such certificate, this
test reports uncertainty rather than declaring the physical system
incapable.
\end{proposition}
\begin{proof}
At a lower face $x_i=b_i$ of the positive box, $x\geq b>0$, so
$A(t)x\geq A_-b$ and monotonicity of saturation gives
\[
        \dot x_i\geq-b_i+\sat((A_-b)_i)-\varepsilon
                    \geq d-\varepsilon>0.
\]
The upper faces retain the same bound because saturation is always at
most one. Odd symmetry proves the negative box invariant. The write and
read calculations are unchanged. Failure to find a sufficient certificate
need not imply failure of the contract, particularly when uncertainty
straddles a threshold, so the final refusal rule follows.
\end{proof}

The finite check in Proposition~\ref{prop:interval-warrant} connects
measured parameter bounds to a proved interface only \emph{conditional}
on the model class and on those bounds remaining valid during service.
It does not turn a finite calibration record into proof of an eternal
transfer law. A known exact rational model permits the complete procedure
of Theorem~\ref{thm:rational-synthesis}; an uncertain physical model may
warrant only a sound positive certificate or an inconclusive result.

\subsection{The two-component construction with explicit bounds}
Take two identical gain-$2$ leaky cells. Initially their hold inputs are grounded. The deployed hold matrix is $A_0=0$. Perform the two generic native link actions
\[
 \mathsf{link}(x_{\rm out},y_{\rm hold}),\qquad
 \mathsf{link}(y_{\rm out},x_{\rm hold}).
\]
The intermediate and final matrices are, up to a choice of coordinate order,
\[
 A_1=\begin{pmatrix}0&0\\2&0\end{pmatrix},\qquad
 A_2=\begin{pmatrix}0&2\\2&0\end{pmatrix}.
\]
No cell state is initialized or transferred by a link. Each link changes which signal drives an input stage. The first link creates a feedforward chain and remains fading; the second closes a cycle. The final Perron root is $2$, whereas $A_0$ and $A_1$ are nilpotent. This is a concrete step at which stable retention becomes possible.

\begin{proposition}[Uniform finite protocol and robustness]\label{fb:prop:two}
For $0\leq\varepsilon<1/2$, the final two-cell device has invariant regions
\[
 B_+=[1/2,1+\varepsilon]^2,
 \qquad B_-=[-1-\varepsilon,-1/2]^2.
\]
Either write command, held for
\begin{equation}\label{fb:eq:two-time}
 T\geq\log\frac{4}{1-2\varepsilon},
\end{equation}
sends every initial state of $[-1-\varepsilon,1+\varepsilon]^2$ into the corresponding region. Thereafter the bit can be retained for every finite hold time and read with any observation error $\nu<1/2$.
\end{proposition}
\begin{proof}
On $B_+$, both nonlinear inputs saturate to $1$, since $2x,2y\geq1$. The hold equations are therefore
\[
 \dot x=1-x+\eta_x,\qquad
 \dot y=1-y+\eta_y.
\]
At either lower face the derivative is at least $1/2-\varepsilon>0$; at an upper face it is at most zero. This proves invariance, and negative symmetry gives $B_-$. The write lower bound $1-\varepsilon-2e^{-T}\geq1/2$ is equivalent to \eqref{fb:eq:two-time}. The coordinate sign is robust to any error less than $1/2$.
\end{proof}
For instance, process disturbance at most $1/4$ permits $T=\log8$ and any read error smaller than $1/2$. With no disturbance, the two boxes contract exponentially to $(1,1)$ and $(-1,-1)$. The saddle at $(0,0)$ remains an equilibrium of the unforced noiseless hold system. The construction does not claim spontaneous or canonical selection of a bit from that symmetric state. The bounded write pulse performs the preparation.

If the actual hold vector field differs uniformly from the specified vector field by at most $\zeta$, and the process disturbance is at most $\varepsilon$, the same invariance proof applies whenever $\varepsilon+\zeta<1/2$. If the write field has the same uniform error bound, replace $\varepsilon$ by $\varepsilon+\zeta$ in the write-time estimate and in the safe outer box. Thus the exact clipping law can be replaced by a uniformly close smooth transfer law on the operating region. This is a conditional implementation guarantee, not evidence that a particular physical circuit satisfies the error bound.

% Optional source module. Uses the theorem/proposition environments of the host.
\subsection{Finite probes can certify a latch without identifying its full transfer law}
Assume each supplied cell has the independently warranted form
\[
 \dot x_i=-x_i+\phi_i(u_i)+\eta_i,
 \qquad |\phi_i(u)|\leq1,\quad |\eta_i|\leq\varepsilon,
\]
where each unknown, time-independent $\phi_i$ is monotone nondecreasing and Lipschitz. This is the same two-cell architecture as before, with the known transfer $\sat(2u)$ replaced by an unknown monotone transfer. The monotonicity, range, leakage law, and noise bounds are model premises, not conclusions extracted from the finite observations below. An external clamp can hold a cell input at a specified real value, and a calibration read has error at most $\nu_{\rm cal}$. Set $M=1+\varepsilon$. The supplied cells begin in $[-M,M]$; no bit value is prescribed. Calibration is performed on the same two physical cells that will subsequently be connected; the model excludes an unreported transfer-law change or hidden lasting damage during a probe. A static setup certificate may be retained, but no copy of a later written data bit is retained outside the deployed device during hold.

\begin{proposition}[Four-probe robust formation certificate]\label{fb:prop:finite-probes}
Choose $b\in(0,1)$ and a finite settling time $\tau$. For each of two cells, clamp its input in turn at $b$ and $-b$ for time $\tau$, and record the endpoint reads $r_i^+$ and $r_i^-$. Put
\[
 \Delta=2e^{-\tau}+\varepsilon+\nu_{\rm cal},\qquad
 c=\min_{i=1,2}\{r_i^+-\Delta,-r_i^--\Delta\}.
\]
If the finite certificate $c>b+\varepsilon$ succeeds, then actually cross-connecting those same two cells produces invariant boxes
\[
 B_+=[b,M]^2,\qquad B_-=[-M,-b]^2
\]
for every transfer-law pair and every disturbance compatible with the model and supplied observations. Clamping both inputs at $b$ or at $-b$ for
\[
 T\geq\log\frac{1+c}{c-\varepsilon-b}
\]
constructs the corresponding box from every initial state in $[-M,M]^2$. Thereafter any read error smaller than $b$ allows indefinite correct bit decoding under hold. Thus the actual records can warrant a reusable capability without reconstructing either full transfer function.
\end{proposition}
\begin{proof}
Under a constant input $u$, variation of constants gives
\[
 |x_i(\tau)-\phi_i(u)|
 \leq (M+1)e^{-\tau}+\varepsilon(1-e^{-\tau})
 =2e^{-\tau}+\varepsilon.
\]
The calibration error adds at most $\nu_{\rm cal}$, so the records imply
$\phi_i(b)\geq c$ and $\phi_i(-b)\leq-c$.

On a lower face of $B_+$, the other coordinate is at least $b$, hence monotonicity gives
$\dot x_i\geq-b+c-\varepsilon>0$. The global range bound gives an inward derivative at the upper faces. The negative-box argument uses the upper bounds at $-b$; oddness is not assumed. This proves robust invariance.

During the positive write, $\phi_i(b)\geq c$, so
\[
 x_i(T)\geq c-\varepsilon-(M+c-\varepsilon)e^{-T}
 =c-\varepsilon-(1+c)e^{-T}\geq b.
\]
The negative write is symmetric at the level of inequalities. The sign decoder is valid with any error smaller than $b$.
\end{proof}

No initialization or reset is required between probes: the bound holds from every state in the invariant ambient box. The probe states can therefore be carried directly through assembly into the proved write protocol. Read rounding is included in $\nu_{\rm cal}$; a certified rational upper bound on $\Delta$ may replace its exact expression. If the strict test fails after the four probes, stop with an inconclusive certificate rather than declaring memory impossible or deploying it as certified.

This is a sound finite certificate, not a complete identification procedure. The test can remain inconclusive for a functioning latch if calibration errors hide the margin. Arbitrarily many finite samples do not establish the monotonicity or global range premise without separate structural warrant. The probes also do not by themselves execute the required cross-connections. Once the certificate has succeeded and the link actions have actually been executed, the previous all-finite-clients reuse argument applies with these boxes and write bounds.

For the concrete gain-$2$ cells $\phi_i(u)=\sat(2u)$, take $b=1/2$, $\varepsilon=1/10$, $\nu_{\rm cal}=1/40$, and $\tau=5$. The conservative rational bound $\Delta=3/20$ is valid because $2e^{-5}+1/10+1/40<3/20$. Every possible four-probe record then gives $c\geq7/10>3/5=b+\varepsilon$. Thus this example succeeds for every admitted calibration disturbance, not merely favorable observations. A write duration $T=3$ is safe because the worst certified ratio is $(1+7/10)/(7/10-1/10-1/2)=17<e^3$. A read error $1/4$ remains safely below the retained-coordinate margin $1/2$. These are finite, fixed, rational settings and known wait times; the
calibration algorithm need not be told that the true transfer is
$\sat(2u)$. The four settling intervals cost twenty cell-time units,
followed by two native link actions. The installation certificate is
fixed before any later data bit arrives, and every later write has the
same three-unit duration.

\subsection{Nondefinability in the whole feedforward baseline}
It would be weak to show merely that a single isolated cell fails. The relevant old language allows any finite feedforward composition of the supplied cells, arbitrary finite fan-out, fixed memoryless Lipschitz interconnection stages, and finite write preparations. It does not allow an output-to-ancestor feedback connection or an external memory that sustains different hold inputs for the two written values. This is the class of already deployed directed-acyclic architectures. Any output-dependent external controller that feeds an ancestor is itself a feedback path and must be counted in the graph; hiding it outside the declared device does not make a feedforward implementation. There is no bound on its finite size or depth, but each proposed implementation has a fixed finite size.

\begin{theorem}[Feedforward fading and robust nondefinability]\label{fb:thm:feedforward}
Every such finite feedforward network forgets its initial state under common future external inputs and common disturbances. More precisely, if the maximum number of dynamic ancestors along a directed path is $d$, differences of states are bounded by $P(t)e^{-t}$ for a finite polynomial $P$ of degree at most $d$, with coefficients depending on the network and initial difference. Consequently no network in this entire baseline implements an indefinitely reusable bit with a fixed positive current-state read margin, even with process disturbance set to zero during hold.
\end{theorem}
\begin{proof}
Topologically order the cells. For a source cell, the two runs have identical input, so its state difference is a constant times $e^{-t}$. For a later cell $i$, let $z_i(t)$ be the absolute state difference. The finite memoryless stages feeding it are Lipschitz, so for fixed finite constants $L_{ij}$ depending on the architecture,
\[
 z_i(t)\leq e^{-t}z_i(0)
       +\int_0^t e^{-(t-s)}\sum_{j<i}L_{ij}z_j(s)\,ds.
\]
If each predecessor is bounded by $P_j(s)e^{-s}$, the right side is
\[
 e^{-t}\left(z_i(0)+\sum_{j<i}L_{ij}\int_0^t P_j(s)\,ds\right),
\]
again polynomial times $e^{-t}$, with degree increasing by at most one along a path. Induction proves the bound and convergence of the complete finite state difference to zero.

Consider any two finite write histories ending in states intended to encode different bits. After both switch to the common hold input, their state difference tends to zero. For any positive observation-error radius, the sets of possible noisy current-state observations eventually intersect. No decoder can return different correct bits on the same possible observation. This contradiction is independent of how the writes were computed, how long they lasted, or which finite feedforward implementation was selected.
\end{proof}

In the purely linear-coupled saturating subfamily, the bound has a compact form. A feedforward coupling matrix $A$ is nilpotent, and coordinatewise comparison gives
\[
 |x(t)-x'(t)|\leq e^{(A-I)t}|x(0)-x'(0)|
 =e^{-t}\sum_{k=0}^{n-1}\frac{t^k A^k}{k!}|x(0)-x'(0)|.
\]
Here the equality uses $A^n=0$, not a heuristic truncation. The larger Lipschitz-network proof above shows that the obstruction is not an accident of one chosen matrix presentation.

A baseline with arbitrarily precise readout, bit-dependent external sustaining inputs, an already persistent controller, or admitted physical feedback is different. The theorem must not be transferred to those baselines. Nor does it exclude retaining a bit for a predetermined finite deadline using sufficiently large feedforward hardware, a tailored precision, or a suitable initial amplitude. The strict result fixes the hardware and error margin while allowing every finite hold duration.

\subsection{The native categorical construction}
\subsubsection{A wiring semantics, not a formal existence assertion}
Use signal interfaces as objects and finite port diagrams of native components as arrows. Primitive boxes include leaky cells, memoryless Lipschitz gains and sums, externally commanded selectors, fan-out, raw state-output ports, and identities. Discontinuous sign decoding and result-dependent command selection belong to the finite macro-client layer, not to an arbitrary instantaneous feedback loop. For each fixed public command, all memoryless stages in the continuous diagram are Lipschitz. A box carries its actual state space, input/output maps, and vector field. Finite parallel composition uses product state spaces. Sequential composition connects outputs to inputs. The allowed diagrams have no algebraic loop made entirely of instantaneous components: every directed cycle must cross a leaky dynamic cell.

This restriction is the native \emph{guarded-feedback} condition. Once the instantaneous acyclic pieces are substituted out, a diagram gives a simultaneous globally Lipschitz ODE on the product of its cell-state spaces. For the circuits used above, this substitution is exactly \eqref{fb:eq:network}. Therefore feedback is implemented by connecting specified ports and solving the resulting actual dynamics. It is not an operation that selects a fixed point of an arbitrary semantic function.

At the syntax level, disjoint union, port connection, renaming, and guarded feedback give the usual symmetric monoidal wiring operations. At the semantic level, two diagrams differing only by reassociation, internal port names, or the order of independent substitutions give the same simultaneous equations up to permutation of state coordinates. Uniqueness of solutions then gives the corresponding equalities of complete trajectories. This proves the coherence needed for the constructions here. We use a guarded or partial trace, not an unqualified total trace on all smooth maps; instantaneous ill-posed feedback is not admitted.

More explicitly, first construct a box with external command input $q$, two feedback inputs $(u,v)$, exposed state outputs $(x,y)$, and two feedback outputs $(y,x)$. Its two cells use $\sat(2u)$ and $\sat(2v)$ during hold and the constant write drive otherwise. Closing its two designated feedback ports gives the two-cell circuit above. Both feedback paths pass through cell dynamics. The exposed outputs remain available for reading. This is a genuine guarded trace of an open, independently specified component diagram; it does not discard the internal state as though an algebraic fixed point had already been reached.

\subsubsection{The baseline and the assembly process are different types of morphism}
The feedforward baseline is the subgrammar of diagrams with a directed-acyclic deployed graph. It is closed under sequential and parallel feedforward substitution, so Theorem~\ref{fb:thm:feedforward} is an invariant of the whole old grammar. A link that closes a cycle is not a morphism of this subgrammar. It is a permitted \emph{formation action} in the larger graph-reconfiguration runtime.

This boundary has a concrete resource meaning. A deployed client owns command and read ports, while the construction process additionally owns unused connection terminals and the authority and actuator to connect them. Initially the material stock has no feedback edge. The two specified generic link operations consume free terminal availability and change the graph. No memory cell with the desired persistent behavior is supplied as an additional primitive. Once connected, the constructed device can be handed to an ordinary client that needs no wiring authority.

One may keep these operations in a two-level category or double category: horizontal morphisms are executions in a fixed architecture, while vertical morphisms are actual graph construction/reconfiguration steps. The present theorems do not presuppose a general double-category formalization or prove a theorem for every contextad. They give concrete native witnesses for both kinds of step. Merely choosing an abstract traced completion without its executable wiring interpretation would not establish physical formation.

\subsubsection{The generated role and reuse are derived}
Let $P=B_-\cup B_+$ and define
\[
 \alpha:P\longrightarrow\{-,+\},\qquad
 \alpha(x)=b\ \text{if }x\in B_b.
\]
This is a generated operational role: its physical representation is the actual state of the newly coupled device, its constructor is the proved write pulse, and its eliminator is the bounded-error read. The invariant-region calculation supplies its warrant rather than assuming a pre-existing native bit interface.

For a fixed autonomous hold law (fixed $A$, or fixed transfer functions
$\phi_i$), and each $t\geq0$, let $H_t\subseteq P\times P$ be the relation of possible endpoints of a hold of duration $t$, over all admitted disturbances. Let $W_b\subseteq E\times P$ be the corresponding duration-$T$ write relation. All endpoint sets are nonempty. In relational semantics the derived equations are
\[
 \alpha H_t=\alpha,\qquad
 \alpha W_b=\operatorname{const}_b.
\]
The first equality expresses robust retention, not merely one chosen successful path. Concatenating disturbances shows $H_{s+t}=H_tH_s$. Splitting a disturbance at the intermediate time proves the reverse inclusion; thus this equality is exact. The noisy read relation followed by the sign decoder is exactly $\alpha$, since the read-error cubes do not cross zero. None of these equalities asserts that the physical states themselves are constant during hold.

\begin{corollary}[All finite typed memory clients]\label{fb:cor:reuse}
Starting with valid states in $P$ (or writing each uninitialized device before its first read), every finite program built from writes of duration at least $T$, arbitrary finite holds, nondestructive reads, result-dependent public branches, and parallel products of independent constructed cells has the same bit-level behavior as the corresponding ideal writable-memory program, for every admitted process and read disturbance. The program may be of arbitrary finite length. Its bounds are not limited to a fixed test suite or a fixed number of uses.
\end{corollary}
\begin{proof}
A write establishes the appropriate invariant region. A hold preserves it, and a read returns its unique label without changing the state. Induction on sequential program formation proves correctness; at a branch the actual read equals the ideal read, so both choose the same continuation. Cartesian products of invariant regions give the parallel statement. The same invariant margins are re-established by every write and are not consumed by a hold or read.
\end{proof}

The compositional law also validates clients with explicit temporary
control state. It does not assert a strict baseline separation for a
client that keeps the written bit externally during a hold. The delayed
comparison below uses the common-hold, no-external-memory contract.

This result concerns explicitly admitted finite clients. Unbounded recursion or liveness properties require their own semantics. It does cover arbitrary finite hold durations, including ones chosen after the hardware was built, and arbitrarily many finite uses. Energy use and physical aging are not free: they need additional operating-lifetime and cost premises in a real device.



\section{A newly executable experiment and coherent later use}\label{sec:experiments}
A writable persistent bit changes which temporal relations can be tested
when no other memory is available. This is a task-level experimental
consequence of the construction, not merely a new name for a state region.

\begin{definition}[Delayed comparison experiment]
A source first returns $b\in\{-1,+1\}$ during a finite preparation phase.
The experimenter may encode it using the admitted write controls, after
which all external bit-dependent control and stored observations are
removed. At a later, previously unspecified finite time, the source
returns $c\in\{-1,+1\}$. A successful experiment returns the bit
$\mathbf1_{b=c}$ using the current plant read and the current input $c$.
The plant is fixed before $b$ arrives. Its hardware size and positive
measurement-error tolerance are fixed independently of the delay.
\end{definition}

\begin{theorem}[Strict gain of an experiment on unchanged boundary bits]
\label{thm:experiment-gain}
A constructed device satisfying Definition~\ref{fb:def:memory} implements
the delayed comparison experiment for every finite delay, every pair of
input bits, and every admitted disturbance and observation error. No
finite feedforward architecture in Theorem~\ref{fb:thm:feedforward}
implements the same experiment with a fixed positive current-state read
margin and no external retained copy of $b$.
\end{theorem}
\begin{proof}
Write $b$, complete the finite write phase, and switch to the common hold
command. At the arrival of $c$, read the stored sign and compare it with
$c$ using a stateless Boolean operation. The proved invariant-region
contract makes the stored sign exactly $b$ for every delay and disturbance.
The comparison is therefore correct.

For a putative feedforward implementation, take its two states after the
finite preparations for $b=+1$ and $b=-1$. With common hold input and zero
process disturbance, their difference tends to zero by
Theorem~\ref{fb:thm:feedforward}. At a sufficiently late common read time,
their positive-radius observation-error sets intersect. Supply the same
second input $c=+1$ in both runs and choose an observation in this
intersection. Every current-read decoder receives identical data, yet
one run requires output one and the other requires output zero. This is
impossible. The delay may depend on the proposed finite implementation,
which is exactly what an unbounded-delay service requirement allows.
\end{proof}

For example, consider two candidate source laws. In one, the second
observation repeats the first; in the other, it reverses it. Let the first bit be fresh and uniformly random in both source laws.
Their late observations then have identical uniform marginals, whereas
the constructed comparison decides the relation exactly on the trial.
This illustration is not a universal statistical identification theorem.
No distributional assumption is needed for the arbitrary-input task
separation just proved.

The new interface can also participate in later construction. In the
reaction model, a manufactured catalyst can be a prerequisite of another
fabrication rule, and the minimal-support and witness results track that
reuse. In the feedback model, a certified memory box can be used as a
component of a finite client, with the abstract bit replacing its hidden
continuous state. Neither substitution makes its material costs or noise
conditions disappear.

\subsection{A general form of the warranted interface law}
The memory example uses invariant relations instead of exact microscopic
reset. This generalizes the point of the state-threaded category without
assuming that a physical device returns to the identical state after every
operation.

Let $X$ be an abstract client-state set, let $P$ be a set of physical
states, and let $R\subseteq P\times X$ relate representations to abstract
states. An abstract operation $f:X\to Y$ is implemented by a nonempty
physical transition relation $F\subseteq P\times Q$ when
\[
 p\mathrel R x, p\mathrel F q
       \quad\Longrightarrow\quad q\mathrel S f(x)
\]
for the chosen output representation relation $S\subseteq Q\times Y$.
Nonemptiness is required at every admitted input representation; it
prevents a vacuous implementation with no executable continuation.
When observations are part of an operation, include their output in the
transition endpoint. Adversarial disturbance choices are quantified over
by this relation.

\begin{proposition}[Composition of guaranteed native realizations]
\label{prop:simulation-composition}
If $F$ implements $f$ relative to $(R,S)$ and $G$ implements $g$ relative
to $(S,T)$, their relational composite implements $g\circ f$ relative to
$(R,T)$. Identity relations implement identities. Independent products
implement products when the native runtime supplies those independent
products and the representation relation is the corresponding product.
If implementation witnesses are retained, their actual executions
concatenate; erasing them yields only the relational statement.
\end{proposition}
\begin{proof}
Given $pRx$ and $p(F;G)r$, choose an intermediate $q$ with $pFq$ and
$qGr$. The first implementation law gives $qS f(x)$ and the second gives
$rT g(f(x))$. Nonemptiness composes by choosing an available first and
then second transition. Identity and product assertions follow directly
from the respective component relations. Concatenation of the chosen
native execution witnesses has the claimed endpoints; no choice of a
new unrelated witness is required after composition.
\end{proof}

For the latch, $R$ is exactly $p\in B_b$. For a catalytic service the
resource stock is restored exactly, which is a special case. The two
native accounts therefore share a compositional interface law while
retaining different physical reuse obligations. This is ordinary
simulation/refinement reasoning. It is used here with constructed
physical regions and native execution programs, rather than as a
substitute for those constructions.

\subsection{What a doctrine change must preserve to retain a gain}
Suppose a comparison translates old native programs into target programs
and respects their observations. That proves preservation of successful
implementations. To preserve a strict gain, one must additionally reflect
the target baseline: any target implementation of the allegedly new task
at the translated old stock must lift to an old baseline implementation,
with the required contract and resources.

\begin{proposition}[Preservation is insufficient without baseline reflection]
\label{prop:service-transport}
Let $T$ be a task, $q_0,q_1$ two source configurations, and let their
target interpretations be $Fq_0,Fq_1$. Assume source service witnesses
translate to valid target witnesses for the same boundary task. Assume
also that every target witness for $T$ at $Fq_0$ lifts to a valid source
witness at $q_0$. Then
\[
 T\in\Cap(q_1)\setminus\Cap(q_0)
 \quad\Longrightarrow\quad
 T\in\Cap'(Fq_1)\setminus\Cap'(Fq_0).
\]
Forward preservation alone does not imply this conclusion, even if the
encoding of configuration labels is injective.
\end{proposition}
\begin{proof}
Translate the positive witness at $q_1$. If a target baseline witness at
$Fq_0$ existed, reflection would lift it to the prohibited source witness.
For failure without reflection, keep all source labels unchanged and add
one target native operation implementing $T$ already at the old stock.
The identity translation preserves all source programs but removes the
separation.
\end{proof}

The proposition identifies the required comparison, not an automatic
method to prove reflection. The finite support test in
Theorem~\ref{syn-thm:transport} supplies such a method for resource
obtainability. It transfers to installed services only when a separate
native theorem identifies exactly which resources implement which
services, as in the retained-catalyst construction following
Theorem~\ref{er:thm:horn}. For the feedback plant, the baseline obstruction
is the fading invariant or the subcritical spectral certificate. A
semantic change that adds a persistent controller, changes the noise
contract, or permits an uncounted feedback edge does not reflect that
baseline.


\section{The resulting constructive claim}\label{sec:combined}
The mathematical conclusion can now be stated without treating a formal
role, a discovered law, and a manufactured device as the same object.

\begin{theorem}[Constructed and reusable capability growth in the stated classes]
\label{thm:combined}
The following are proved construction procedures under independently
specified native models.
\begin{enumerate}
\item In a finite positively graded executable reaction system, the
algorithm of Theorem~\ref{er:thm:synthesis} computes obtainable stable
stocks and exactly stock-preserving services, with native manufacture
and use words and independently checkable obstruction certificates.
Corollary~\ref{er:cor:doctrine} computes the complete bounded client
capability category, including costs. The retained-catalyst subclass
realizes Horn formation with exact finite feedstock accounting and an
exact service interpretation.
\item In the cooperative saturating cell model, a rational irreducible
candidate architecture can be classified by finite positive or negative
certificates. A successful candidate is compiled into native links,
finite writes, invariant hold regions, and a positive-margin read.
Finite candidate lists admit complete selection within that list.
\item For the larger class of unknown monotone leaky transfer cells,
the four-probe certificate provides a sound finite experimental route
to constructing the same type of reusable bit, conditional on its
stated structural and error promises. A failed certificate is not a
negative decision about the plant.
\item The acquired interfaces compose by their native reset or invariant
relations. The latch enables the delayed comparison experiment, while
the entire stated finite feedforward baseline cannot implement it.
The formation comparison and its conservative transport criterion are
exact in the retained-resource subclass.
\end{enumerate}
These are gains in installed or ready-service capability, with explicitly
charged construction and external inputs. They do not add an operation
outside the complete construction-and-use runtime's extensional closure.
\end{theorem}
\begin{proof}
The first item is the finite reachability and shortest-path construction,
with the elementary service interpretation following
Theorem~\ref{er:thm:horn}. The second is
Theorems~\ref{fb:thm:threshold} and~\ref{thm:rational-synthesis}; the
positive certificates supply every invariant, write, and read parameter.
The third is Proposition~\ref{fb:prop:finite-probes}. Native service words
concatenate by Corollary~\ref{er:cor:reuse}; invariant-region clients
compose by Corollary~\ref{fb:cor:reuse} and
Proposition~\ref{prop:simulation-composition}. The experiment separation
is Theorem~\ref{thm:experiment-gain}; formation and its transport are
Theorems~\ref{syn-thm:coupling} and~\ref{syn-thm:transport}. Finally,
Theorem~\ref{thm:closure} applies to every full native construction and
compiled continuation. No part of the proof requires treating a new
semantic name as a supplied physical primitive.
\end{proof}

This theorem is a synthesis of specific constructions, not a claim that
one generic closure theorem solves invention. Its difficult obligations
are discharged differently in different native classes. The reaction
model has finite exact search and exact reset. The feedback model has
continuous states, robust invariant sets, and a whole-baseline fading
obstruction. Their common categorical content is state-threaded native
composition and guaranteed interface simulation. Their physical
assumptions are not interchangeable.

\section{Connection to the earlier native programme}\label{sec:project}
The previous \emph{Operational Invention from Interaction} constructed
image roles of certified regular interactions, acquired frames, and used
retained references to change accessible operations. It proved a
resource-relative capability result while leaving a stronger endogenous
construction target open. A frame or reference can be valuable without
being a newly fabricated dynamical mechanism. The present latch starts
instead with components that individually lose the required information
and produces a robust retention capacity by changing their interaction.
The catalyst instance likewise starts without the active service resource
and returns an elementary fabrication trace that creates it.

The distinction is not that the new artifacts violate a fixed substrate.
They satisfy that substrate's laws throughout. The substantive additions
are explicit manufacture, the comparison to its pre-installation service,
a native use protocol, and a proof covering subsequent experiments.
The closure boundary from the earlier note remains valid and is used
rather than bypassed.

The earlier \emph{Native Operations Across Semantic Doctrines} separated
\[
       \text{source grammar}\longrightarrow
       \text{native realizations}\longrightarrow
       \text{semantic operations}.
\]
It supplied native-recognition and reflection results under finite probe
and coverage conditions, and showed how witness-erasing doctrine changes
can create operations without compatible source lifts. Here the middle
term additionally depends on an installed resource state. The maps are
therefore indexed by that state, and a manufacturing arrow changes the
index. Proposition~\ref{er:prop:reset} shows why an arbitrary such arrow
does not automatically induce a conservative service functor.

There are two positive recognition mechanisms in the present note.
Minimal-support antichains recognize what a persistent formation grammar
can obtain. Rational invariant boxes and finite calibration inequalities
recognize a robust dynamical interface. Neither is the circular test
``accept exactly when an implementation exists.'' Both can be computed
or checked from independently specified primitive data and finite
certificates. Their scope is still their own native classes.

The programme's proof-relevant native setting retains more than an
extensional relation: execution occurrences, endpoint data, ownership,
costs, and comparison witnesses may all matter. In the reaction system,
firing words are retained beside the reachability or cost abstraction.
In the dynamic system, the native wiring actions and command execution
are retained beside the bit-level simulation. Passing to a preorder,
cost matrix, or relation explicitly forgets information. None of these
abstractions reconstructs discarded histories.

A lift into the full contextad or relative-module setting would require
its specified endpoint transport, units, multiplication, pure maps, and
comparison cells, with the appropriate coherence and resource laws.
The guarded wiring semantics here provides native equations and
compositional executions for these particular circuits. It is not a
proof that every abstract contextad, traced completion, or semantic
operation has a physical realization. Existing contextad results remain
at their stated strength; they must not be renamed as new results of
this note~\cite{CM2026}.

\section{Relations to established work}\label{sec:prior}
The target has close antecedents. The appropriate claim is a proved
construction for the stated native classes with explicit comparisons.
The exact combined formulation has not received an exhaustive priority
search, and historical novelty is not asserted.

\subsection{Executable objects and resources that act on further objects}
Fontana and Buss's algorithmic chemistry is a direct generative
antecedent~\cite{Fontana1992,FontanaBuss1994}. Executable expressions
interact to produce further expressions, with entities functioning as
operations on other entities. The programme distinguishes the underlying
combinatory universe from the organizations actually available in a
current stock. Accordingly, making operations into objects, synthesizing
new executable combinations, and letting those combinations participate
in later interactions are not new ideas of this note. The present
finite reaction theorem adds an explicit bounded synthesis and
stock-preserving service certificate for its model; it does not displace
algorithmic chemistry's broader constructive programme.

Coecke, Fritz, and Spekkens already construct resource theories of states,
parallel-combinable processes, and finite stocks of processes combined
through admissible circuit contexts~\cite{CFS2016}. The state-threaded
service category here is compatible with that established perspective.
Its extra obligations are the native fabrication witness, the exact
returned-stock or invariant-region law, and the baseline obstruction.
No resource-theoretic category alone guarantees that a one-use process
can be reused indefinitely.

Deutsch's constructor theory makes repeatability central: a constructor
performs a task while retaining the ability to perform it again, not
necessarily the identical microscopic state~\cite{Deutsch2013}. That
is closely reflected in the latch's invariant regions. The accounting
requirement is also the same: energy supplies, substrates, and waste
cannot be hidden. The present results are effective procedures for
restricted finite or dynamical classes, not a replacement for the
possible/impossible task programme.

\subsection{Dynamical memory and categorical wiring}
Positive feedback, bistability, sequential circuits, and persistent
memory are established engineering and mathematical ideas. Hopfield's
continuous-response networks give a direct example of collective
memory generated by coupled graded-response units~\cite{Hopfield1984}.
The present two-cell circuit is a mathematical latch model, not a new
memory invention. The proofs use standard Perron--Frobenius theory,
comparison estimates, invariant rectangles, and variation of constants.
The spectral threshold is stated for the precise saturating cooperative
class, and its historical novelty is not established.

Vagner, Spivak, and Lerman give an established symmetric monoidal
wiring-diagram account of open dynamical systems~\cite{VSL2015}.
The native guarded wiring used here instantiates that kind of semantics
with explicit cell laws and executable connections. Its operational
benefit is not the word ``feedback'' or a formal trace. It is that the
specified port connections yield a well-posed ODE whose invariant regions
support the claimed service. An arbitrary categorical feedback operator
without that implementation argument would not establish the theorem.

The strict comparison also has a precise scope. Acyclic networks have a
fading invariant under the admitted common hold inputs. A network with
hidden feedback, an external persistent controller, or infinitely precise
sensing is a different baseline. The result does not manufacture
persistent memory out of a source already silently allowed to retain the
bit by one of those means.

\subsection{Learning reusable operations and experimental interfaces}
DreamCoder already alternates learning reusable executable abstractions
with solving further synthesis problems~\cite{DreamCoder}. Higher-order
inductive logic programming likewise combines abstraction and predicate
invention with learning advantages~\cite{CropperMuggleton2016}.
These systems show why an executable library may be genuinely useful
without exceeding its interpreter's extensional closure. The present
construction instead measures an installed physical service and proves a
robust task-level obstruction for a declared earlier architecture.
Neither comparison invalidates the other.

Konidaris's skill--symbol loop and the theorem-backed construction of
planning symbols from skills already connect operational acquisition
with representation formation~\cite{Konidaris2016,KKL2018}.
CALF and categorical automata learning reconstruct behavioral objects
under explicit access and termination hypotheses~\cite{CALF,FunLstar}.
CoalA additionally refines component alphabets and learns a decomposition
in its synchronizing-system setting~\cite{CoalA2025}. It would therefore
be incorrect to claim that all earlier active learning fixes every
component interface in advance. The distinction made here is between
learning a behavioral interface and constructing a new physical mechanism
that implements it, with a subsequent reuse theorem.

Tool design and use have also been connected in robotics.
Liu and collaborators learn task-conditioned tool designs and
corresponding controllers and demonstrate their approach on a physical
robot~\cite{Liu2023}. A general slogan that learning creates a usable
tool is therefore insufficient as a contribution. The present note
provides mathematical construction and impossibility certificates for
its own models; it does not claim a robot performed its fabrication or
that its abstract runtime solves real-world tool design.

\subsection{Changing interfaces and theory regimes}
Spivak's polynomial trees explicitly model interfaces reshaped by
interaction, with morphisms, tensor structure, and monoidal
closure~\cite{Spivak2026}. The changing-interface idea and its categorical
organization are direct prior work. An additional theorem is needed to
infer or manufacture one of those interfaces in a physical substrate.
The present note does not claim an equivalence between polynomial trees
and its reaction or feedback models.

Recent theory-revision systems also go beyond fitting parameters within
one fixed model class. Ma and collaborators jointly choose structural
edits and diagnostic experiments while explicitly recognizing limits of
their edit library~\cite{Ma2026}. Wang and Buehler describe verified
representation-regime transitions and categorical transport of earlier
artifacts~\cite{WangBuehler2026}. These are close comparators to the broader
programme, not evidence that categorical theory change is unoccupied.
The native construction, source-warrant, service-preservation, and
baseline-reflection obligations above specify what must be compared at
the theorem level.

Taylor's distinctions among exploratory, expansive, and transformational
open-endedness emphasize that novelty is relative to a possibility
space and its level of description~\cite{Taylor2019}. The present results
establish finite reachable manufacture and reusable services. They do
not establish open-ended evolution, autonomous discovery of arbitrary
native laws, or an indefinitely expanding universe of operational kinds.

\section{Verification and remaining limits}\label{sec:verification}
All principal mathematical claims in this source have proofs. The source
is deliberately uncompiled; no PDF layout or TeX compilation result is
claimed.

The companion \texttt{EndogenousGrowth.lean} was checked with Lean 4.35.0
and mathlib revision
\texttt{56a005855ea17e0b0958f73ec801ae140b0a1679}.
It formalizes the persistent finite-premise closure core, leastness and
invariant obstructions, finite construction plans and their interpretation
using independently supplied native rule implementations, finite
saturation, forward compiler composition, reflection-based strict-gap
transport, and a joint-formation example. Its finite saturation reaches
a fixed point by the number of capability labels. Counterexamples show
that even an injective forward compiler need not preserve a strict gain.
The checked source uses no unproved placeholders or \texttt{native\_decide};
its reported axioms are the standard
\texttt{propext}, \texttt{Classical.choice}, and \texttt{Quot.sound}.

This formalization does not verify the whole manuscript. In particular,
it does not formalize mass-consuming reaction execution, shortest-path
optimality, the minimal-support algorithm, dynamical systems, spectral
criteria, finite calibration, resource costs, physical realization, or
historical priority. The finite reaction and closure regression checks
provide independent executable examples and counterexamples, not proofs
of the general physical premises. Numerical trajectories are not used
as evidence for an indefinite robustness theorem.

The exact limits are therefore substantive. The elementary reactions must
be executable as specified. The cells must obey the warranted leakage,
range, regularity, and error model throughout the claimed operating
period. Readout must be nondestructive or have separately accounted
effects. Powered operation consumes energy; indefinite ideal holding is
a mathematical statement about all finite times in a stationary model,
not a promise of infinite physical lifetime. Hardware fabrication,
source-law acquisition, and transfer to a real material system require
additional evidence.

Within those limits, the stronger target has a concrete positive answer.
Interaction can construct an artifact with a derived reusable interface,
change the installed experiments that are possible, and preserve the
interface under subsequent native composition. The answer is neither a
new name for an old semantic map nor a claim to escape fixed physical
laws. It is an explicitly constructed and checked capability change,
with the exact comparison and its obstructions left visible.


\begin{thebibliography}{99}
\bibitem{Felleisen1991}
M. Felleisen, \emph{On the expressive power of programming languages},
Science of Computer Programming 17 (1991), 35--75.
\url{https://www.cs.tufts.edu/comp/150FP/archive/matthias-felleisen/expressive-as-published.pdf}.

\bibitem{Fontana1992}
W. Fontana, \emph{Algorithmic Chemistry}, 1992.
\url{https://sites.santafe.edu/~walter/Papers/alchemy92.pdf}.

\bibitem{FontanaBuss1994}
W. Fontana and L. W. Buss,
\emph{What would be conserved if ``the tape were played twice''?},
Proceedings of the National Academy of Sciences 91 (1994), 757--761.
\url{https://pmc.ncbi.nlm.nih.gov/articles/PMC43028/}.

\bibitem{CFS2016}
B. Coecke, T. Fritz, and R. W. Spekkens,
\emph{A mathematical theory of resources},
Information and Computation 250 (2016), 59--86.
\url{https://arxiv.org/abs/1409.5531}.

\bibitem{Deutsch2013}
D. Deutsch, \emph{Constructor Theory}, Synthese 190 (2013), 4331--4359.
\url{https://arxiv.org/abs/1210.7439}.

\bibitem{Hopfield1984}
J. J. Hopfield, \emph{Neurons with graded response have collective
computational properties like those of two-state neurons},
Proceedings of the National Academy of Sciences 81(10) (1984), 3088--3092.
\url{https://doi.org/10.1073/pnas.81.10.3088}.

\bibitem{VSL2015}
D. Vagner, D. I. Spivak, and E. Lerman,
\emph{Algebras of open dynamical systems on the operad of wiring diagrams},
Theory and Applications of Categories 30 (2015), 1793--1822.
\url{https://arxiv.org/abs/1408.1598}.

\bibitem{DreamCoder}
K. Ellis et al., \emph{DreamCoder: Growing generalizable, interpretable
knowledge with wake-sleep Bayesian program learning},
arXiv:2006.08381 (2020).
\url{https://arxiv.org/abs/2006.08381}.
The related PLDI 2021 paper is titled
\emph{DreamCoder: Bootstrapping Inductive Program Synthesis with
Wake-Sleep Library Learning}.

\bibitem{CropperMuggleton2016}
A. Cropper and S. H. Muggleton,
\emph{Learning Higher-Order Logic Programs through Abstraction and
Invention}, IJCAI 2016, 1418--1423.
\url{https://ijcai.org/proceedings/2016/204}.

\bibitem{Konidaris2016}
G. Konidaris, \emph{Constructing Abstraction Hierarchies Using a
Skill-Symbol Loop}, IJCAI 2016.
\url{https://arxiv.org/abs/1509.07582}.

\bibitem{KKL2018}
G. Konidaris, L. P. Kaelbling, and T. Lozano-P\'erez,
\emph{From Skills to Symbols: Learning Symbolic Representations for
Abstract High-Level Planning},
Journal of Artificial Intelligence Research 61 (2018), 215--289.
\url{https://doi.org/10.1613/jair.5575}.

\bibitem{CALF}
B. van Heerdt, M. Sammartino, and A. Silva,
\emph{CALF: Categorical Automata Learning Framework}, CSL 2017.
\url{https://arxiv.org/abs/1704.05676}.

\bibitem{FunLstar}
T. Colcombet, D. Petri\c{s}an, and R. Stabile,
\emph{Learning Automata and Transducers: A Categorical Approach},
CSL 2021.
\url{https://arxiv.org/abs/2010.13675}.

\bibitem{CoalA2025}
L. Henry, M. R. Mousavi, T. Neele, and M. Sammartino,
\emph{Compositional Active Learning of Synchronizing Systems Through
Automated Alphabet Refinement}, CONCUR 2025.
\url{https://drops.dagstuhl.de/entities/document/10.4230/LIPIcs.CONCUR.2025.20}.

\bibitem{Liu2023}
Z. Liu, S. Tian, M. Guo, K. Liu, and J. Wu,
\emph{Learning to Design and Use Tools for Robotic Manipulation},
Proceedings of Machine Learning Research 229 (2023), 887--905.
\url{https://proceedings.mlr.press/v229/liu23b.html}.

\bibitem{Spivak2026}
D. I. Spivak, \emph{Interactions that reshape the interfaces of the
interacting parties}, arXiv:2602.17917v1 (2026).
\url{https://arxiv.org/abs/2602.17917v1}.

\bibitem{Ma2026}
S. Ma et al., \emph{Learning to Outgrow a Theory: Experimental Discovery
Beyond the Initial Hypothesis Space}, arXiv:2610.07627 (2026).
\url{https://arxiv.org/abs/2610.07627}.

\bibitem{WangBuehler2026}
F. Y. Wang and M. J. Buehler,
\emph{Self-Revising Discovery Systems for Science: A Categorical
Framework for Agentic Artificial Intelligence},
arXiv:2606.01444v2 (2026).
\url{https://arxiv.org/abs/2606.01444}.

\bibitem{Taylor2019}
T. Taylor, \emph{Evolutionary Innovations and Where to Find Them:
Routes to Open-Ended Evolution in Natural and Artificial Systems},
Artificial Life 25(2) (2019), 207--224.
\url{https://arxiv.org/abs/1806.01883}.

\bibitem{CM2026}
M. Capucci and D. J. Myers,
\emph{Contextful Arrows via Wreath Products},
arXiv:2410.21889v2, revised 3 November 2026.
\url{https://arxiv.org/abs/2410.21889v2}.

\end{thebibliography}
\end{document}
